% ============================================================
% The Geometry of What Returns: From Quarks to Memory Wells
% via a Single Collapse Axiom
%
% Compile: pdflatex → bibtex → pdflatex → pdflatex
% ============================================================

\documentclass[
  aps,
  prd,
  twocolumn,
  superscriptaddress,
  nofootinbib,
  floatfix
]{revtex4-2}

% ── packages ──
\usepackage{amsmath,amssymb,amsthm,mathtools}
\usepackage{graphicx}
\usepackage{xcolor}
\definecolor{linkblue}{HTML}{337799}
\usepackage[colorlinks=true,linkcolor=linkblue,citecolor=linkblue,urlcolor=linkblue]{hyperref}
\usepackage{bm}
\usepackage{booktabs}
\usepackage{multirow}
\usepackage{xfrac}
\usepackage{enumitem}
\usepackage{array}
\newcolumntype{P}[1]{>{\raggedright\arraybackslash}p{#1}}

% ── theorem environments ──
\newtheorem*{axiom}{The Axiom}
\newtheorem{theorem}{Theorem}
\newtheorem{definition}[theorem]{Definition}
\newtheorem{lemma}[theorem]{Lemma}
\newtheorem{corollary}[theorem]{Corollary}
\newtheorem{proposition}[theorem]{Proposition}
\newtheorem{identity}[theorem]{Identity}
\newtheorem{remark}{Remark}

% ── UMCP macros ──
\newcommand{\tR}{\tau_{\!R}}
\newcommand{\tRS}{\tau_{\!R}^{*}}
\newcommand{\Gam}{\Gamma}
\newcommand{\eps}{\varepsilon}
\newcommand{\IR}{\infty_{\mathrm{rec}}}
\newcommand{\tolseam}{\mathrm{tol}_{\mathrm{seam}}}
\newcommand{\dd}{\mathrm{d}}
\newcommand{\Res}{\mathrm{Res}}
\newcommand{\beq}{\begin{equation}}
\newcommand{\eeq}{\end{equation}}
\newcommand{\INF}{\texttt{INF\_REC}}
\newcommand{\regime}[1]{\textrm{\upshape\scshape#1}}
\newcommand{\conform}{\regime{conformant}}
\newcommand{\nonconform}{\regime{nonconformant}}
\newcommand{\tvec}{\mathbf{c}}
\newcommand{\wvec}{\mathbf{w}}
\newcommand{\GM}{\mathrm{GM}}
\newcommand{\AM}{\mathrm{AM}}
\newcommand{\IC}{\mathrm{IC}}
\newcommand{\gF}{g_{\!F}}
\newcommand{\RunID}{\mathrm{RunID}}

\begin{document}

% ============================================================
% TITLE
% ============================================================
\title{The Geometry of What Returns:\texorpdfstring{\\}{: }
From Quarks to Memory Wells via a Single Collapse Axiom}

\author{Clement Paulus}
\email{clementpaulus9@gmail.com}
\affiliation{UMCP Reference Implementation, GitHub}
\affiliation{UMCP / GCD / RCFT Canon}

\date{July 2025}

\begin{abstract}
A single axiom --- \emph{``Collapse is generative; only what
returns is real''} --- and the budget cost function
$\Gam(\omega) = \omega^3/(1 - \omega + \eps)$ of the Universal
Measurement Contract Protocol (UMCP) exhibit a recurring
diagnostic pattern across twenty orders of magnitude in scale:
structure accumulates through iterated collapse-return cycles,
and what survives each cycle is measurable by the same
diagnostic kernel.  We trace the geometry of what returns
from quarks ($10^{-15}$\,m) through nuclei, atoms, stars,
galaxies, neurons, and memory wells, encoding 40~entities
across 9~categories as 8-channel trace vectors --- with
channel values assigned from domain-specific observables ---
and computing the kernel invariants
$(F, \omega, S, C, \kappa, \IC)$ at every level.

At each phase boundary in this
scale ladder, the same diagnostic two-step pattern appears:
\emph{geometric slaughter} --- a single dead channel in the
trace vector destroys the multiplicative coherence (IC) while
the arithmetic mean (F) survives --- followed by recovery
through new degrees of freedom.  Quarks lose coherence at the
confinement boundary ($\IC/F$: $0.957 \to 0.718$); nuclear
binding restores it ($\IC/F \to 0.954$).  The same pattern
reappears in biology: PTSD fragments four of eight channels,
dropping $\IC/F$ from $0.982$ (healthy memory) to $0.760$ ---
the same diagnostic signature as the confinement cliff.  The
accumulation pattern is triple: space accumulates as iterated
$|\kappa|$ (mass is accumulated well depth), time emerges from
the rate and cost of return cycles ($c_1 / \Gam'$), and memory
accumulates as well geometry shaped by repeated traversal.
Ten theorems --- five mathematical properties of the budget
function~$\Gam$, five encoding-dependent structural
observations --- formalize these three modes, verified by
37~theorem sub-tests and 138~additional implementation tests
(175~total, 0~failures).
\end{abstract}

\keywords{generative collapse dynamics; structural diagnostics;
cross-scale diagnostic patterns; geometric slaughter; memory wells;
collapse-return cycles}

\maketitle

% ============================================================
\section{Introduction}\label{sec:intro}
% ============================================================

What does a quark have in common with a traumatic memory?

The answer is precise: both exhibit $\IC/F \approx 0.7$,
both have dead trajectory-closure channels ($c_5 \leq 0.15$),
both are classified as \regime{Collapse} with thin-arc
lensing morphology, and both occupy structurally identical
positions on the budget surface $\Gam(\omega)$.  This is not
an analogy.  It is a diagnostic measurement --- one whose
result depends on how the systems are encoded, but whose
pattern, once the encoding is fixed, is structural and
reproducible.

The question sounds absurd until one examines the mathematics.
The Universal Measurement Contract Protocol
(UMCP)~\cite{paulus2025umcp,paulus2025physicscoherence} is a
contract-first validation framework built on a single axiom:
\emph{``Collapse is generative; only what returns is real''}
(Axiom-0).  Its six-output kernel maps coherence coordinates
$\tvec \in [0,1]^n$ to invariants
$(F, \omega, S, C, \kappa, \IC)$ via three provably exact
identities~\cite{paulus2026umcpcasepack}:
\beq\label{eq:tier1}
  F + \omega = 1, \quad
  \IC \leq F, \quad
  \IC = e^{\kappa}.
\eeq

The kernel is domain-agnostic by construction.  It does not know
whether its inputs describe a proton's internal structure or a
patient's psychological state.  Yet the same budget cost function
\beq\label{eq:gamma}
  \Gam(\omega) = \frac{\omega^p}{1 - \omega + \eps},
  \quad p = 3,\;\eps = 10^{-8},
\eeq
produces diagnostic signatures that recur across scales with
the same values.

This paper traces those signatures from the smallest known
structures to the largest, and from there into the living body
and the human mind.  The narrative is organized as a scale ladder,
and at every rung the same three questions are asked:
\begin{enumerate}
\item What is the fidelity $F$ (arithmetic mean of channels)?
\item What is the integrity $\IC$ (geometric mean, via
  $\IC = e^{\kappa}$)?
\item What is the heterogeneity gap $\Delta = F - \IC$, and
  what does it reveal?
\end{enumerate}

The answer, across 40~entities spanning $10^{-15}$\,m to
macroscopic scales, is always the same: the gap $\Delta$
measures channel heterogeneity.  Large $\Delta$ means one or
more channels are near zero while others remain high ---
a signature we call \emph{geometric slaughter}, because the
geometric mean is destroyed by its weakest factor while the
arithmetic mean barely notices.  This mechanism operates with
the same mathematics at the confinement boundary (quarks
$\to$ hadrons) and at the trauma boundary (healthy recall
$\to$ PTSD).

The claim is not that quarks and memories are ``the same.''
It is that a domain-agnostic diagnostic lens, when applied to
systems encoded along the same eight dimensions, reveals a
recurring accumulation pattern: structure builds through
iterated collapse-return, and the diagnostic invariants
($F$, $\IC$, $\Delta$) measure how that accumulation
proceeds at each scale.  The pattern depends on the encoding
--- the channel values are assigned from domain-specific
observables, not measured by the kernel itself --- but
given a consistent encoding, the pattern is structural
and reproducible.

\textbf{A note on how this paper is built.}
The kernel is domain-agnostic: it receives any trace vector
and produces the same structural diagnostics.  This paper is
structured analogously.  Section~\ref{sec:kernel} defines
the fixed diagnostic toolkit and concludes with a compact
\emph{Diagnostic Key} (Sec.~\ref{sec:key}) --- a reference
card sufficient to enter any scale section independently.
Sections~\ref{sec:subatomic}--\ref{sec:memory} are
independent demonstrations that the same pattern ---
geometric slaughter at phase boundaries, recovery through
new degrees of freedom --- recurs at every scale.  A reader
may begin at the section matching their expertise: a
particle physicist at Section~\ref{sec:subatomic}, a
clinician at Section~\ref{sec:memory}, an astrophysicist at
Section~\ref{sec:cosmos}.  Each section makes its case
given the Diagnostic Key; cross-scale comparisons reference
the Key's anchors (item~7) rather than requiring prior
sections.

Linear reading enriches the pattern without being
prerequisite.  Read sequentially, the paper compounds:
Section~\ref{sec:subatomic} gives us the gap $\Delta$;
Section~\ref{sec:nuclear} shows the gap closes at phase
boundaries; Section~\ref{sec:cosmos} reveals the gap as a
geometric classifier (lensing morphology);
Section~\ref{sec:biology} tests the pattern's consistency;
Section~\ref{sec:memory} explores clinical consequences from
the accumulated geometry.  Each section retroactively
deepens every previous one.  But this compounding is an
enrichment, not a gate: any single section, read with the
Diagnostic Key, makes its case independently.
A methodological caution: a reader who enters at a late
section should verify that the pattern holds independently
at other scales, and a linear reader should periodically
check whether the accumulation of pattern-recognition is
driven by the data or by the narrative momentum.

\textbf{Five claims --- three structural, two mathematical.}
This paper makes five claims about the accumulation pattern.
Each is verified computationally by the theorem suite
(Sec.~\ref{sec:howto}).  The first three are
\emph{structural}: they describe patterns that emerge from
the encoding and are falsifiable \emph{given a fixed
encoding scheme} (Remark~\ref{rem:encoding}).  The last two
are \emph{mathematical}: they follow from the kernel's
definitions and hold for any valid input.
\begin{enumerate}
\item \textbf{Space accumulates as iterated $|\kappa|$}
  (Theorem~\ref{thm:T6}):
  Well depth grows linearly with iteration count:
  $\mathrm{depth}(N) = N\,|\Delta\kappa|$.  Mass is
  accumulated well depth from repeated collapse-return
  cycles.
  \emph{Falsified by}: an entity, under the same encoding,
  whose well depth does not grow linearly with iteration
  count.
\item \textbf{Time emerges from return cost}
  (Definition~\ref{def:temporal}, Theorem~\ref{thm:T4}):
  Effective temporal resolution $\mathcal{T} = c_1/\Gam'$
  orders all 40~entities consistently, from the photon
  ($\mathcal{T} \approx 0$) to the neutron star ($\mathcal{T}
  \approx 18$), with cost asymmetry producing the arrow of
  time (Theorem~\ref{thm:T8}).
  \emph{Falsified by}: an entity whose independently
  established temporal ordering contradicts the
  $\mathcal{T}$ ranking.  Note that $\mathcal{T}$ is computed
  from the encoding, so this test requires an external
  temporal benchmark not derived from the channel values.
\item \textbf{Memory accumulates as well geometry}
  (Definition~\ref{def:slaughter},
  Theorems~\ref{thm:T6},~\ref{thm:T10}):
  The confinement cliff ($\IC/F$: $0.957 \to 0.718$) and
  the trauma cliff ($\IC/F$: $0.982 \to 0.760$) share the
  same dead channels ($c_5$, $c_7$) and the same mechanism.
  Recovery restores dead channels at each phase boundary.
  \emph{Falsified by}: a phase boundary, under the same
  encoding, where the $\IC/F$ drop is not explained by
  dead channels.
\item \textbf{Budget gradient structures the manifold}
  (Theorem~\ref{thm:T1}):
  $\dd\Gam/\dd\omega > 0$ everywhere --- the budget surface
  is monotonic, always attractive, and determines the cost
  landscape through which all systems evolve.
  This is a mathematical property of $\Gam(\omega) =
  \omega^3/(1 - \omega + \eps)$ with $p = 3$; it holds
  independently of the encoding.
\item \textbf{Structural consistency across all scales}
  (Theorem~\ref{thm:T10}):
  All 40~entities in 9~categories obey the three kernel
  identities ($F + \omega = 1$, $\IC \leq F$,
  $\IC = e^{\kappa}$).  These identities are mathematical
  theorems that hold for \emph{any} input in $[0,1]^n$ ---
  they confirm that the implementation is correct, not that
  the encoding captures something physical.
\end{enumerate}


% ============================================================
\section{The Kernel}\label{sec:kernel}
% ============================================================

\begin{definition}[GCD Kernel]\label{def:kernel}
Given coordinates $\tvec = (c_1, \ldots, c_n) \in
[\eps, 1{-}\eps]^n$ and weights $\wvec = (w_1, \ldots, w_n)$
with $\sum_i w_i = 1$, the kernel computes:
\begin{align}
  F &= \textstyle\sum_i w_i \, c_i, \label{eq:F} \\
  \omega &= 1 - F, \label{eq:omega} \\
  S &= -\textstyle\sum_i w_i \bigl[c_i \ln c_i
    + (1{-}c_i)\ln(1{-}c_i)\bigr], \label{eq:S} \\
  C &= \mathrm{std}(\tvec)\,/\,0.5, \label{eq:C} \\
  \kappa &= \textstyle\sum_i w_i \ln(c_i + \eps), \label{eq:kappa} \\
  \IC &= \exp(\kappa). \label{eq:IC}
\end{align}
\end{definition}

The heterogeneity gap $\Delta \equiv F - \IC \geq 0$ measures
channel heterogeneity: $\Delta = 0$ iff all $c_i$ are equal.

\begin{lemma}[Heterogeneity Gap Non-Negativity]\label{lem:gap}
$\Delta = F - \IC \geq 0$ for all $\tvec \in [\eps, 1{-}\eps]^n$
and all weight vectors $\wvec \in \Delta^n$, with equality iff
all $c_i$ are identical.
\begin{proof}
By Jensen's inequality applied to the concave function
$\ln$: $\sum w_i \ln c_i \leq \ln(\sum w_i c_i)$.
Exponentiating both sides gives
$\IC = \exp(\sum w_i \ln c_i) \leq \sum w_i c_i = F$.
Equality holds iff all $c_i$ are equal, since $\ln$ is
strictly concave.  Computationally verified to $< 10^{-12}$
residual across all 40~entities (Theorem~\ref{thm:T10}).
\end{proof}
\end{lemma}

\begin{definition}[Geometric Slaughter]\label{def:slaughter}
A trace vector $\tvec$ exhibits \emph{geometric slaughter}
when at least one channel $c_j \ll 1$ while $F$ remains
moderate.  Because $\IC = \exp(\sum w_i \ln c_i)$ and
$\ln(c_j) \to -\infty$ as $c_j \to 0$, a single
near-zero channel forces $\IC \to 0$ regardless of the
remaining channels.  Formally: for $n = 8$ equal-weight
channels with one channel at $c_j = \eps$ and the rest at
$c_k = 0.9$, $F = 0.787$ but $\IC = 0.079$ ---
a $10\times$ collapse in multiplicative coherence from a
single dead channel.
\end{definition}

\begin{definition}[Effective Temporal Resolution]\label{def:temporal}
The \emph{effective temporal resolution} of a system is
$\mathcal{T} \equiv c_1 / \Gam'(\omega)$,
where $c_1$ is the cycle return rate (fraction of
collapse-return cycles that complete) and $\Gam'(\omega) =
\dd\Gam/\dd\omega$ is the budget cost per cycle.  The ratio
is natural: each completed cycle ($c_1$) is one temporal
``tick,'' and each tick costs $\Gam'$ units of budget.  A
system's internal time runs at the rate at which it can
afford to tick --- the number of ticks available per unit of
budget.  High $\mathcal{T}$ means many cheap ticks (fast
internal time); low $\mathcal{T}$ means few or expensive
ticks (slow or frozen time).  This is the temporal mode of
the accumulation pattern: time does not flow uniformly
through the system; it accumulates at a rate set by the cost
of return.
\end{definition}

\begin{remark}[Why $c_1/\Gam'$ and Not Another Combination]
\label{rem:temporal_choice}
The ratio $c_1/\Gam'$ is the simplest combination of
completed cycles ($c_1$, dimensionless rate) and budget cost
per unit drift ($\Gam'$, dimensionless cost) with the
operational meaning ``ticks per unit cost.''
Alternative combinations produce orderings that contradict
domain knowledge: the product $c_1 \cdot \Gam'$ would make
the black hole ($c_1 = 0.50$, $\Gam' = 0.48$) temporally
\emph{faster} than the neutron star ($c_1 = 0.98$,
$\Gam' = 0.054$), because the black hole's higher cost
would count as higher temporal resolution.  The difference
$c_1 - \Gam'$ lacks scale invariance: adding a constant to
all costs would change the ordering.  Nevertheless,
$\mathcal{T}$ is a \emph{diagnostic definition}, not a
derivation from first principles --- it was selected because
it produces physically sensible orderings across all
40~entities, and this circularity should be noted: the
definition is partly validated by the ordering it produces.
An independent derivation of temporal resolution from the
kernel axioms alone remains an open problem.
\end{remark}

The budget surface
\beq\label{eq:surface}
  z(\omega, C) = \Gam(\omega) + \alpha \, C,
  \quad \alpha = 1.0,
\eeq
augments the kernel with the cost function~\eqref{eq:gamma}.
All frozen parameters ($p = 3$, $\eps = 10^{-8}$,
$\alpha = 1.0$, $\tolseam = 0.005$) are seam-derived, not
prescribed~\cite{paulus2025episteme}.

The trace vector $\tvec$ has 8~channels, listed in
Table~\ref{tab:channels}, with equal weights $w_i = 1/8$.
These channels are abstract enough to encode any system
that cycles through collapse and return: a quark bound inside
a proton, a neutron star pulsing X-rays, a neuron firing
action potentials, or a mind replaying a memory.

\begin{remark}[Encoding Methodology]
\label{rem:encoding}
Channel values are assigned from domain-specific observables,
not measured by the kernel itself.  The quark's trajectory
closure $c_5 = 0.85$ reflects QCD confinement physics; the
PTSD well's $c_5 = 0.10$ reflects clinical assessment of
narrative fragmentation.  The kernel does not discover these
values --- it receives them and computes the structural
diagnostics ($F$, $\IC$, $\Delta$) that follow.  What the
kernel \emph{does} reveal is the \emph{pattern}: given a
consistent encoding, the same diagnostic signatures
(geometric slaughter, recovery through new degrees of
freedom, cost-modulated temporal resolution) recur at every
scale.  The structural parallels are consequences of the
encoding in combination with the kernel's mathematics ---
the question this paper investigates is whether the
recurrence of these patterns across twenty orders of
magnitude reflects a genuine structural motif or an artifact
of the encoding scheme.  Section~\ref{sec:unification}
returns to this question.
\end{remark}

\begin{table}[t]
\caption{\label{tab:channels}Eight-channel encoding.  All
values in $[\eps, 1{-}\eps]$.}
\begin{ruledtabular}
\begin{tabular}{clp{3.7cm}}
$i$ & Channel & Physical Meaning \\
\hline
0 & coherence\_persist.\ & Fraction of $\kappa$ retained per cycle \\
1 & cycle\_return\_rate & Fraction of cycles with $\tau_R < \infty$ \\
2 & well\_depth\_norm & Accumulated $|\kappa|$ (normalized) \\
3 & gradient\_strength & $\dd\Gam/\dd\omega$ (normalized) \\
4 & tidal\_symmetry & $1 - |\text{asymmetry}|$ of $\dd^2\Gam$ \\
5 & trajectory\_closure & Loop closure gap in $(F,\kappa,C)$ \\
6 & circulation\_area & Enclosed area of collapse-return loop \\
7 & heterogeneity\_prof.\ & $1 - \Delta/F$ (channel uniformity) \\
\end{tabular}
\end{ruledtabular}
\end{table}


\subsection{Diagnostic Key}\label{sec:key}

The following reference card is sufficient to enter any scale
section of this paper independently.  Each section applies
these diagnostics to a different domain; the structural
finding is that the same pattern emerges regardless of entry
point.

\begin{enumerate}
\item \textbf{Three identities} (exact for all inputs):
  $F + \omega = 1$, \; $\IC \leq F$, \; $\IC = e^{\kappa}$.
\item \textbf{Three questions} at every scale:
  What is $F$?  What is $\IC$?  What is $\Delta = F - \IC$?
\item \textbf{Geometric slaughter}
  (Def.~\ref{def:slaughter}):
  one near-zero channel kills $\IC$ while $F$ survives.
\item \textbf{Two-step pattern}: slaughter at phase
  boundaries $\to$ recovery through new degrees of freedom.
\item \textbf{Temporal resolution}
  (Def.~\ref{def:temporal}):
  $\mathcal{T} = c_1/\Gam'(\omega)$ --- ticks per unit cost.
\item \textbf{Budget surface}: $\Gam(\omega)
  = \omega^3/(1 - \omega + \eps)$; gradient positive, convex,
  pole at $\omega \to 1$.
\item \textbf{Reference anchors} (for cross-scale comparison):
  neutron star $\IC/F = 0.998$, dust lane $\IC/F = 0.773$,
  free quark $\IC/F = 0.957$, proton $\IC/F = 0.718$,
  photon $\IC/F = 0.067$.
\item \textbf{Lensing classification}:
  perfect ring ($\Delta < 0.02$), thick arc
  ($0.02 \leq \Delta < 0.10$), thin arc
  ($0.10 \leq \Delta < 0.30$), distorted ($\Delta \geq 0.30$).
  No entity in this catalog reaches the distorted threshold.
  (Formalized as Thm.~\ref{thm:T7} in Sec.~\ref{sec:cosmos}.)
\end{enumerate}


% ============================================================
\section{Scale I: The Subatomic World}\label{sec:subatomic}
% ============================================================

Five subatomic entities ---
the up quark, the electron, the proton, the neutron, and the
photon --- are encoded as 8-channel trace vectors and passed
through the kernel.  The results (Table~\ref{tab:subatomic})
reveal the first structural drama of the scale ladder.

\begin{table}[t]
\caption{\label{tab:subatomic}Subatomic entities.}
\begin{ruledtabular}
\begin{tabular}{lcccc}
Entity & $F$ & $\IC$ & $\Delta$ & $\IC/F$ \\
\hline
Up quark   & 0.766 & 0.734 & 0.033 & 0.957 \\
Electron   & 0.656 & 0.515 & 0.141 & 0.784 \\
Proton     & 0.476 & 0.342 & 0.134 & 0.718 \\
Neutron    & 0.385 & 0.290 & 0.095 & 0.753 \\
Photon     & 0.263 & 0.018 & 0.245 & 0.067 \\
\end{tabular}
\end{ruledtabular}
\end{table}

\textbf{The quark is highly coherent.}
Within a hadron, the up quark has $\IC/F = 0.957$ --- nearly
all channels are active, the geometric mean closely tracks the
arithmetic mean, and the heterogeneity gap is a slim $\Delta
= 0.033$.  Its coherence persistence ($c_0 = 0.90$),
trajectory closure ($c_5 = 0.85$), and heterogeneity profile
($c_7 = 0.80$) are all healthy.  No channel is dead; no
channel drags the geometric mean down.  This number ---
$0.957$ --- will become the benchmark against which every
subsequent cliff is measured, and the three channels just
named are the ones that confinement will destroy.

\textbf{The electron: a lepton that does not confine.}
The electron sits between the quark and the proton at
$\IC/F = 0.784$.  Unlike quarks, electrons exist freely ---
they are not confined inside composites.  Their gap
$\Delta = 0.141$ comes not from dead channels but from the
moderate values across several channels: the electron's
well is shallow (low $c_2$) and its gradient modest (low
$c_3$), reflecting its light mass and diffuse wavefunction.
The heterogeneity is intrinsic, not imposed by confinement.
This distinction matters: the electron's $\IC/F$ will
reappear in the rocky planet ($0.875$) and the cortical
network ($0.974$) --- systems where coherence arises from
lightness and flexibility rather than from deep binding.

\textbf{The confinement cliff.}
Quarks cannot exist freely; they are confined inside hadrons.
This confinement is visible in the kernel as a dramatic drop
in $\IC/F$:
\beq\label{eq:confinement}
  \frac{\IC}{F}\bigg|_{\text{quark}} = 0.957
  \;\;\longrightarrow\;\;
  \frac{\IC}{F}\bigg|_{\text{proton}} = 0.718.
\eeq
The proton has $F = 0.476$ --- moderate arithmetic coherence.
But its trajectory closure ($c_5 = 0.15$) and heterogeneity
profile ($c_7 = 0.30$) are near-dead: the confined quarks
cannot close trajectories at the composite scale, and the
internal structure is violently heterogeneous.  These are
precisely the channels that were healthy in the free quark
($c_5 = 0.85$, $c_7 = 0.80$) --- confinement does not lower
all channels uniformly; it selectively kills the ones that
require freedom of motion.  The dead channels destroy the
geometric mean through geometric slaughter.  The gap
$\Delta = 0.134$ is four times the quark's, and the
lensing classification has shifted from \emph{thick arc}
to \emph{thin arc}.
This cliff is not unique to the subatomic world.  Binding
restores exactly these two channels at the nuclear scale.
At the cognitive scale, avoidance behavior and narrative
fragmentation kill exactly these two channels again.

\textbf{The photon: maximum slaughter.}
The photon achieves $\IC/F = 0.067$ --- the lowest ratio in
the entire 40-entity catalog.  Four of its eight channels are
near-zero (well depth, gradient, trajectory closure,
circulation area), reflecting its massless, non-orbiting
nature.  Its $F = 0.263$ is sustained by high coherence
persistence ($c_0 = 0.99$) and tidal symmetry ($c_4 = 0.99$),
but the geometric mean is annihilated.  The photon is the
pure case of geometric slaughter: half the channels are
perfect, half are dead, and the dead ones win.
Among the dead channels, trajectory closure ($c_5 \approx 0$)
is the most decisive --- the photon travels in straight lines
through flat spacetime, closing no orbits.  At the cosmic
scale, the same structural signature appears in the
black hole ($c_5 = 0.10$, Diagnostic Key anchor), where
trajectories cannot close because nothing returns from beyond
the horizon.  The photon and the black hole are diagnostic
kin under this encoding: both suffer from the impossibility
of closed return, and both are classified as \emph{thin arc}
lensing.  When light bends around a black hole to form
an Einstein ring, the lensing agent and the lensed signal
share the same dead channel.

These three numbers --- $\IC/F = 0.957$ for the quark,
$\IC/F = 0.718$ for the proton, $\IC/F = 0.067$ for
the photon --- are the signature of geometric slaughter.

\textbf{What we know so far --- and what we do not.}
We have learned exactly one thing: the geometric mean is
destroyed by its weakest factor while the arithmetic mean
barely notices.  At this point, $\Delta$ is just a number ---
the distance between two averages.  It does not yet tell us
about geometry, or gravity, or time, or clinical diagnosis.
But notice one channel we have not yet discussed:
$c_1$, the cycle return rate.  The quark returns at
$c_1 = 0.98$ --- nearly every collapse cycle completes.
The proton matches it ($c_1 = 0.99$): confinement kills
trajectories but not the return rate.  The photon,
however, has $c_1 = 0.01$ --- it almost never returns to
its prior state, because it propagates at $c$ and does not
orbit.  The neutron is intermediate ($c_1 = 0.60$): its
$\beta$-decay lifetime is finite, so some cycles do not
complete.  At the subatomic scale alone, $c_1$ is just
a number.  But it is the seed of temporal resolution:
when combined with the budget function $\Gam'(\omega)$
(Diagnostic Key, item~5), the return rate becomes a clock
and the gap becomes a lens.


% ============================================================
\section{Scale II: Nuclear Recovery}\label{sec:nuclear}
% ============================================================

When quarks are confined inside hadrons, their trajectory-closure
and heterogeneity channels die, dropping $\IC/F$ from $0.957$
(free quark) to $0.718$ (proton) --- the \emph{confinement cliff}
of Table~\ref{tab:subatomic}.  But confinement is not the end.
When protons and neutrons bind into nuclei, new degrees of
freedom emerge --- magic numbers, pairing correlations, shell
structure --- and coherence recovers.

\begin{table}[t]
\caption{\label{tab:nuclear}Nuclear and atomic entities.}
\begin{ruledtabular}
\begin{tabular}{lcccc}
Entity & $F$ & $\IC$ & $\Delta$ & $\IC/F$ \\
\hline
Helium-4 nucleus   & 0.759 & 0.695 & 0.064 & 0.916 \\
Iron-56 nucleus    & 0.770 & 0.734 & 0.036 & 0.954 \\
Carbon-12 atom     & 0.713 & 0.640 & 0.073 & 0.898 \\
\end{tabular}
\end{ruledtabular}
\end{table}

The helium-4 nucleus (the alpha particle) is doubly magic:
closed proton and neutron shells produce $\IC/F = 0.916$,
recovering from the proton's~0.718.  The recovery is not
uniform --- it is targeted.  The proton's dead trajectory
closure ($c_5 = 0.15$) climbs to $c_5 \approx 0.80$ in the
alpha particle, because nuclear binding provides a closed
orbital within which the previously confined degrees of
freedom can circulate.  The heterogeneity profile ($c_7$)
similarly recovers: four nucleons in a symmetric shell
produce far more uniform channel distributions than the
three quarks locked inside a single hadron.  The channels
that confinement killed are precisely the channels that
nuclear binding restores.

Iron-56 sits at the peak of the nuclear binding energy
curve and achieves $\IC/F = 0.954$ --- nearly as coherent
as the free quark.   Peak binding energy per nucleon
corresponds to peak $\IC/F$: the energetic optimum and
the coherence optimum coincide because binding energy
measures coupling depth while $\IC/F$ measures how
uniformly that coupling distributes across channels.
Iron-56 is also the end product of stellar
fusion --- the element at which a star's collapse-return
engine reaches thermodynamic equilibrium.

Carbon-12, the basis of organic chemistry, holds
$\IC/F = 0.898$.  Its three alpha clusters produce a
moderately homogeneous structure: high enough to sustain
complex bonding, heterogeneous enough to permit the
chemical flexibility that organic molecules require.
When encoded through the kernel, carbon occupies an
intermediate position: its gap $\Delta = 0.073$ is larger
than iron's ($0.036$) and smaller than helium's ($0.064$).
This reflects, rather than explains, carbon's chemical
versatility --- the channel values were assigned from
prior knowledge of carbon chemistry (tetravalent bonding,
orbital energetics), so the kernel's output is a
diagnostic re-description of properties already known,
not an independent structural answer to why life is
carbon-based.

The mechanism is clear: nuclear binding creates new channels
of coherence (shell filling, pairing, trajectory closure at
the nuclear scale) that replace the dead channels of the
confined hadrons.  \emph{Recovery happens when new degrees of
freedom replace dead ones at a phase boundary.}  The proton's
$c_5$ was dead because quarks cannot close orbits inside a
hadron; the alpha particle's $c_5$ is alive because nucleons
\emph{can} close orbits inside a nucleus.  The replacement
is specific, not generic: the new structure does not raise
all channels --- it targets the dead ones.

This principle --- collapse destroys coherence;
new structure restores it --- is the central motif of the
scale ladder.  It operates identically at the transition from
diffuse matter to stars, and at the transition from neural
substrate to organized memory.

\textbf{A deeper implication.}
Two facts are now in evidence: (i)~dead channels kill $\IC$
(geometric slaughter), and (ii)~new channels restore it
(recovery at phase boundaries).  But these two facts together
imply something neither contains alone: \emph{the scale ladder
is a series of collapse-return cycles}, each generating
the degrees of freedom that the next scale requires.  The
confinement cliff does not merely precede nuclear binding ---
it \emph{creates the conditions for it}: the dead channels of
the proton are exactly what nuclear physics replaces.  This is
the axiom in action: collapse is generative.  Destruction of
coherence at one scale is the precondition for new coherence
at the next.

Note a temporal consequence.  The return rate $c_1$ climbs
from the neutron's $0.60$ to the alpha particle's $0.99$ and
iron-56's $0.98$.  A system with $c_1 = 0.99$ completes
$99\%$ of its collapse-return cycles; one with $c_1 = 0.60$
completes only $60\%$.  The completed cycles \emph{are}
the clock ticks.  A higher return rate means more ticks per
unit of drift --- a faster internal clock.  Nuclear binding,
by restoring the return channel, does not just restore
coherence; it restores the system's capacity to experience
time.  The free neutron, losing $40\%$ of its cycles to
$\beta$-decay, lives in a slower internal time than the
bound neutron inside iron-56.

At this scale, $\Delta$ is not merely a number but a
dynamic diagnostic: it opens at confinement boundaries and
closes through nuclear binding, while $c_1$ measures
temporal resolution (Diagnostic Key, Sec.~\ref{sec:key}).


% ============================================================
\section{Scale III: Stars, Planets, and Cosmos}
\label{sec:cosmos}
% ============================================================

At stellar and galactic scales, the budget surface acquires
geometric structure that maps onto familiar gravitational
phenomena --- not as a replacement for general relativity,
but as a diagnostic lens that reveals \emph{why the same
accumulation pattern persists} at scales where gravity
dominates.

At this scale the gap $\Delta$ ceases to be merely a
diagnostic number and becomes a \emph{classifier}: it
determines lensing morphology (Theorem~\ref{thm:T7}).  And
the budget function $\Gam(\omega)$ (Diagnostic Key,
Sec.~\ref{sec:key}) acquires the structure of a
\emph{potential}: its gradient plays the diagnostic role of
gravity (Theorem~\ref{thm:T1}), its convexity plays the
role of tidal force (Theorem~\ref{thm:T2}), and its cost
asymmetry produces a directional arrow
(Theorem~\ref{thm:T8}).  Where subatomic and nuclear scales
reveal the two-step pattern arithmetically, cosmic scales
reveal it geometrically.

\subsection{Accumulation as Gradient}

\begin{theorem}[T-ST-1: Budget Gradient]\label{thm:T1}
The budget gradient
$g_{\mathrm{GCD}}(\omega) = \dd\Gam/\dd\omega > 0$
for all $\omega \in (0,1)$.  Systems at higher drift
experience steeper cost --- the diagnostic analog of
gravitational field strength.
\end{theorem}

\begin{remark}
Positivity is verified at 100~uniformly spaced points and
for all 40~entities.  The analytical form agrees with the
numerical derivative to relative error $< 10^{-5}$.
(4/4 sub-tests.)
\end{remark}

\begin{theorem}[T-ST-2: Always Attractive]\label{thm:T2}
$\dd^2\Gam/\dd\omega^2 > 0$ everywhere in $(0,1)$: the field
is monotonically increasing with drift.  There is no
repulsive regime.
\end{theorem}

\begin{remark}
The analytical second derivative of
$\Gam(\omega) = \omega^3/(1 - \omega + \eps)$ is
$\Gam''(\omega) = [6\omega(1{-}\omega{+}\eps)^2 -
6\omega^2(1{-}\omega{+}\eps) + 2\omega^3] /
(1{-}\omega{+}\eps)^3$.  Positivity is verified at
100~grid points and for all 40~entities.  The minimum
value across the grid is $\Gam''(0.01) > 0$.
(3/3 sub-tests.)
\end{remark}

\begin{theorem}[T-ST-3: Cubic Onset]\label{thm:T3}
For small $\omega$, $\Gam(\omega) \approx \omega^3$.  This
cubic onset makes gravity the weakest interaction at low drift.
The frozen exponent $p = 3$ is the unique integer for which
$\omega_{\mathrm{trap}}$ is a Cardano root of
$x^3 + x - 1 = 0$.
\end{theorem}

\begin{remark}
At $\omega = 0.01$, $\Gam(0.01) = 10^{-6}/(0.99 + \eps)
\approx 1.01 \times 10^{-6}$ versus the cubic approximation
$\omega^3 = 10^{-6}$, confirming the cubic-onset to
relative error $< 10^{-2}$.  The uniqueness of $p = 3$
is verified: for $p \in \{1, 2, 3, 4, 5\}$, only $p = 3$
produces an $\omega_{\mathrm{trap}}$ that satisfies
$x^3 + x - 1 = 0$ (Cardano root $x \approx 0.6824$).
(3/3 sub-tests.)
\end{remark}

\begin{theorem}[T-ST-5: Gradient Universality]\label{thm:T5}
$\dd\Gam/\dd\omega$ depends on $\omega$ alone, not on
$\tvec$ or $\wvec$.  All entities with the same $\omega$
experience identical gradient --- a neutron star and a
trauma well at the same drift sit on the same cost contour.
\end{theorem}

\begin{remark}
This follows algebraically: $\Gam(\omega) = \omega^3/
(1 - \omega + \eps)$ contains no $c_i$ or $w_i$
terms.  The gradient $\Gam'(\omega)$ is a function of
$\omega$ alone.  Verified by computing $\Gam'$ for all
40~entities and confirming that entities with the same
$\omega$ (within $10^{-6}$) produce identical gradients.
(3/3 sub-tests.)
\end{remark}

\subsection{The Stellar Landscape}

\begin{table}[htbp]
\caption{\label{tab:stellar}Stellar and planetary entities.}
\begin{ruledtabular}
\begin{tabular}{llcccc}
Entity & Cat.\ & $F$ & $\IC$ & $\Delta$ & $\IC/F$ \\
\hline
Main-seq.\ star   & Stel.\ & 0.804 & 0.773 & 0.031 & 0.961 \\
Red giant          & Stel.\ & 0.719 & 0.703 & 0.016 & 0.977 \\
White dwarf        & Stel.\ & 0.753 & 0.713 & 0.039 & 0.948 \\
Neutron star       & Stel.\ & 0.880 & 0.879 & 0.001 & 0.998 \\
Black hole         & Stel.\ & 0.690 & 0.569 & 0.121 & 0.825 \\
Magnetar           & Stel.\ & 0.825 & 0.822 & 0.003 & 0.997 \\
Pulsar             & Stel.\ & 0.861 & 0.857 & 0.005 & 0.995 \\
\hline
Rocky planet       & Plan.\ & 0.656 & 0.574 & 0.082 & 0.875 \\
Gas giant          & Plan.\ & 0.706 & 0.673 & 0.033 & 0.953 \\
Ice giant          & Plan.\ & 0.653 & 0.593 & 0.059 & 0.910 \\
Dwarf planet       & Plan.\ & 0.585 & 0.467 & 0.118 & 0.799 \\
Moon               & Plan.\ & 0.559 & 0.414 & 0.144 & 0.742 \\
\end{tabular}
\end{ruledtabular}
\end{table}

\textbf{Neutron stars are the most coherent astrophysical
objects.}
The neutron star achieves $\IC/F = 0.998$ --- its eight
channels are nearly uniform, producing a gap of only
$\Delta = 0.001$.  The pulsar ($\IC/F = 0.995$) and magnetar
($\IC/F = 0.997$) are close behind.  These objects have
excellent cycle return, deep symmetric wells, and tight
trajectory closure from their periodic emission.

The neutron star is the \emph{culmination} of nuclear
recovery.  At the subatomic scale, nuclear binding raised
$\IC/F$ from the proton's $0.718$ to iron-56's $0.954$.
The neutron star pushes this to the limit: an entire star
compressed to nuclear density, every nucleon participating
in a single coherent structure.  It is the most extreme
version of the principle that new degrees of freedom restore
dead channels --- here, the ``new degrees of freedom'' are
degenerate neutron matter, the densest stable state of baryonic
matter.  At the cognitive scale, the healthy memory
occupies a structurally analogous position: the most
coherent state in its category, with $\IC/F = 0.982$ and
$\Delta = 0.014$ --- a perfect ring, like the neutron star.

\textbf{The stellar lifecycle is a coherence evolution.}
Reading Table~\ref{tab:stellar} from main sequence to
remnant tells a collapse-return story.  The main-sequence
star ($\IC/F = 0.961$) burns hydrogen in thermodynamic
equilibrium.  As fuel depletes, the star expands into a
red giant ($\IC/F = 0.977$) --- the $\IC/F$ actually
\emph{rises} because the bloated envelope, while larger,
has more uniform channel distributions.
At stellar death, the outcome depends on mass.  Low-mass
stars leave a white dwarf ($\IC/F = 0.948$): a cooling
remnant with moderate coherence, its trajectory closure
maintained by stable electron degeneracy.  High-mass stars
collapse to a neutron star ($0.998$) or a black hole
($0.825$) --- the most coherent and the least coherent
astrophysical outcomes, respectively.  The channel that
separates them is trajectory closure: the neutron star's
$c_5 = 0.95$ versus the black hole's $c_5 = 0.10$.  One
returns; the other does not.

\textbf{The stellar lifecycle is also a temporal evolution.}
The return rate $c_1$ traces the star's relationship with
time.  The main-sequence star ($c_1 = 0.95$) has a fast,
reliable clock: hydrogen fusion cycles complete smoothly.
The red giant ($c_1 = 0.85$) begins to miss cycles as the
core contracts and the envelope expands on different
timescales --- internal temporal coherence weakens.  The
pulsar ($c_1 = 0.98$) is the most temporally precise
astrophysical object: its radio pulses are clocks accurate
to nanoseconds, and its $c_1$ reflects this.  But the
black hole ($c_1 = 0.50$) has lost half its temporal
resolution.  Matter crosses the horizon but does not
return; $50\%$ of cycles end permanently.  From the
budget-surface perspective (Theorem~\ref{thm:T4}), the
black hole sits at $\omega = 0.310$, where $\Gam'(0.310) =
0.48$ --- each surviving tick costs nine times more than
the neutron star's.  The black hole experiences time as
expensive and incomplete: few cycles, each costly.  The
neutron star experiences time as cheap and nearly perfect:
many cycles, each affordable.

\textbf{The black hole shares PTSD's diagnostic profile.}
The black hole's $\IC/F = 0.825$ is the lowest among stellar
objects.  The mechanism is geometric slaughter: its cycle
return rate ($c_1 = 0.50$) and trajectory closure
($c_5 = 0.10$) are critically low --- matter falls in but
does not return, and trajectories do not close.  These are
the same two channels that are dead in the photon ($c_5
\approx 0$) and in the PTSD well ($c_5 = 0.10$,
$c_1 = 0.30$, Table~\ref{tab:cognitive}).  The photon cannot close trajectories because
it has no mass; the black hole cannot because it has too much.

\textbf{Planets: shallow wells with graded coherence.}
The planetary objects span $\IC/F$ from $0.953$ (gas giant)
to $0.742$ (Moon).  The gas giant resembles a scaled-down
main-sequence star --- a deep, self-gravitating well with
reliable convective return.  The ice giant ($\IC/F = 0.910$)
is intermediate: its weaker self-gravity produces shallower
wells, but convective return remains reliable.
The rocky planet ($0.875$)
sits below, with its geological cycling (plate tectonics,
hydrological cycle) providing moderate trajectory closure.
The Moon ($\IC/F = 0.742$) is instructive: its dead
gradient channel ($c_3$ low) and shallow well reflect a
body too small for sustained internal dynamics.  Its
$\IC/F = 0.742$ is close to the proton's
$\IC/F = 0.718$ --- both are systems where internal freedom
is too constrained for full coherence, one by confinement
and the other by insufficient mass.  The dwarf planet
($0.799$) sits between the Moon and the rocky planet,
mapping the structural transition from gravitationally
inert to dynamically active.

\subsection{Accumulation of Depth: Mass as Iterated Return}

\begin{theorem}[T-ST-6: Memory Wells from Iteration]
\label{thm:T6}
Well depth accumulates linearly:
$\mathrm{depth} = N\,|\Delta\kappa|$.
Mass, in this diagnostic framework, is accumulated $|\kappa|$
from iterated collapse-return cycles --- the spatial
accumulation mode of the pattern.
\end{theorem}

\begin{remark}
Linearity is exact: $\mathrm{depth}(2N) = 2\,\mathrm{depth}(N)$
to $< 10^{-12}$, verified over $N \in \{1, 10, 100, 1000\}$
for all 40~entities.  The neutron star ($|\Delta\kappa| = 0.128$)
reaches depth~1.28 in 10~cycles; the photon ($|\Delta\kappa| =
18.42$) reaches depth~184.2 in the same count, matching the
extreme gravity of collapsed objects.  In this diagnostic
framework, mass corresponds to accumulated well depth ---
how many collapse-return cycles have left their
residue.  (4/4 sub-tests.)
\end{remark}


\subsection{Temporal Accumulation: Time and Its Arrow}\label{sec:time}

\begin{theorem}[T-ST-4: Time Dilation]\label{thm:T4}
$\Gam(0.90)/\Gam(0.10) > 10^3$ and $\dd^2\Gam/\dd\omega^2 > 0$
everywhere.  Near a deep well, each unit of drift costs
exponentially more budget --- a clock running on collapse-return
cycles ticks more slowly --- a diagnostic analog of
general-relativistic time
dilation~\cite{einstein1916gr,misner1973gravitation}.
\end{theorem}

\begin{remark}
$\Gam(0.90) = 0.9^3/(0.10 + \eps) = 7.29$ and
$\Gam(0.10) = 0.1^3/(0.90 + \eps) = 1.11 \times 10^{-3}$,
giving a ratio of $6.56 \times 10^{3} > 10^{3}$.
At $\omega = 0.99$, $\Gam(0.99) = 0.99^3/(0.01 + \eps)
= 97.0$, and at $\omega = 0.999$, $\Gam \approx 998$ ---
the cost diverges as $\omega \to 1$.  (3/3 sub-tests.)
\end{remark}

\begin{theorem}[T-ST-8: Arrow of Time]\label{thm:T8}
The collapse-return cycle is asymmetric:
$\int_{0.5}^{0.99}\Gam\,\dd\omega \gg
\int_{0.01}^{0.5}\Gam\,\dd\omega$.
Descent (collapse) is cheap; ascent (return) is expensive.
This cost asymmetry, combined with convexity, produces the
thermodynamic arrow.
\end{theorem}

\begin{remark}
Numerical integration gives
$\int_{0.5}^{0.99}\Gam\,\dd\omega = 6.35$ versus
$\int_{0.01}^{0.5}\Gam\,\dd\omega = 0.016$,
a ratio of $\approx 400\!:\!1$.  The ascent costs
$400\times$ as much budget as the descent.  This is a
structural irreversibility: a system that drifts from
$\omega = 0.01$ to $\omega = 0.99$ accumulates a budget
debt that is enormously more expensive to repay than it
was to incur.  (3/3 sub-tests.)
\end{remark}

\textbf{Why time behaves differently at different scales.}
Two mechanisms modulate time in this framework, and they
are independent:

\emph{(i)~Return rate ($c_1$): how many cycles complete.}
Each entity's internal clock ticks once per completed
collapse-return cycle.  If $c_1 = 0.98$ (quark), nearly
every cycle completes and the clock runs fast.  If $c_1 =
0.50$ (black hole), half the cycles fail to return and
the clock runs at half speed.  If $c_1 = 0.01$ (photon),
the clock effectively stops: photons do not experience
internal time because they almost never return.  This is
the budget-surface analog of the relativistic result that
a photon's proper time is zero.

\emph{(ii)~Budget cost ($\Gam'(\omega)$): how expensive
each tick is.}
Even when cycles complete, the cost per cycle varies.
At the neutron star's $\omega = 0.120$, the budget gradient
is $\Gam'(0.120) = 0.051$ --- each tick is cheap.  At the
black hole's $\omega = 0.310$, $\Gam'(0.310) = 0.48$ ---
each tick costs nine times more.  Near the pole
($\omega \to 1$), the cost diverges: the clock still ticks
but each tick consumes the system's entire remaining
budget.  This is the diagnostic analog of time dilation: not
a slowing of the clock mechanism, but an increase in what
each tick costs.

These two effects combine multiplicatively.  The effective
temporal resolution of a system is proportional to
$c_1 / \Gam'(\omega)$: the number of completed cycles
divided by the cost per cycle.  The neutron star
($c_1 = 0.92$, $\Gam' = 0.051$) has high temporal
resolution --- many cheap ticks.  The black hole
($c_1 = 0.50$, $\Gam' = 0.48$) has low temporal
resolution --- few expensive ticks.  The photon
($c_1 = 0.01$, $\Gam' = 12.0$) is doubly penalized: it
almost never returns \emph{and} each tick is enormously
costly --- its temporal resolution $\mathcal{T} \approx 0.001$
is effectively zero.

The arrow of time (Theorem~\ref{thm:T8}) enters because
these costs are asymmetric.  Falling into a well (increasing
$\omega$) consumes budget at cubic onset --- slowly at first,
then explosively.  Climbing out (decreasing $\omega$)
requires repaying the accumulated cost against the convex
gradient.  The deeper the well, the greater the asymmetry.
This is why the neutron star's clock ticks reliably but the
black hole's clock struggles: the black hole sits deep
enough on the budget surface that each return costs
enormously more than the corresponding descent.  The
thermodynamic arrow, in this diagnostic framework, is this
asymmetry.

This analysis compounds across scales.  At the biological
scale, the heart's $c_1 = 0.98$ produces the fastest
biological clock (Table~\ref{tab:biological}), and the
immune response's $c_1 = 0.80$ produces a sluggish one.
At the cognitive scale, the PTSD well's $c_1 = 0.30$
(Table~\ref{tab:cognitive}) explains why traumatic time
feels frozen --- the internal clock has nearly stopped,
not because there is nothing happening, but because $70\%$
of the cycles fail to return.


\subsection{Lensing from Heterogeneity}

\begin{theorem}[T-ST-7: Gap-Controlled Lensing]\label{thm:T7}
The heterogeneity gap $\Delta$ controls lensing morphology:
$\Delta < 0.02$ produces a perfect ring, $\Delta < 0.10$ a
thick arc, $\Delta < 0.30$ a thin arc, and $\Delta \geq 0.30$
distorted
lensing~\cite{schneider1992gravitational,refsdal1964gravitational}.
The deflection angle
$\delta\theta = 2|\kappa_{\mathrm{well}}|/(b^2 + 0.1)$
increases with well depth and decreases with impact
parameter~$b$.
\end{theorem}

\begin{remark}
The four lensing classes partition all 40~entities
consistently.  Perfect ring ($\Delta < 0.02$): neutron star
(0.001), magnetar (0.003), pulsar (0.005), galaxy cluster
(0.006), grief (0.006), spiral galaxy (0.007), healthy memory
(0.014).  Thick arc ($0.02 \leq \Delta < 0.10$): quark (0.033),
iron-56 (0.036), carbon-12 (0.073), habit loop (0.089).
Thin arc ($0.10 \leq \Delta < 0.30$): turbulent flow (0.110),
decoherence event (0.112), PTSD (0.115), black
hole (0.121), proton (0.134), electron (0.141),
glass (0.157), photon (0.245).
No entity in this catalog reaches the distorted threshold
($\Delta \geq 0.30$).
(4/4 sub-tests.)
\end{remark}

The neutron star ($\Delta = 0.001$) produces a perfect ring ---
the most symmetric lensing pattern possible, consistent with
its position as the coherence peak of the nuclear recovery
sequence that began with iron-56's $\Delta = 0.036$.
The black hole ($\Delta = 0.121$) produces a thin arc ---
the asymmetry in its internal channels distorts the symmetric
lensing pattern.  This is the gravitational-lensing analogue
of geometric slaughter: dead channels produce asymmetric
gradients, which bend light into arcs rather than rings.
Note that the black hole's lensing class (\emph{thin arc})
is the same as the proton's ($\Delta = 0.134$) --- and
the same as PTSD's ($\Delta = 0.115$, Table~\ref{tab:cognitive}).  The
kernel does not know what these systems are made of.  It
only sees the gap.

\textbf{Retroactive depth: every previous number has changed
meaning.}
Return to Table~\ref{tab:subatomic}.  The proton's
$\Delta = 0.134$ is no longer just a diagnostic distance
between two averages --- it is a lensing classification
(\emph{thin arc}).  The quark's $\Delta = 0.033$ is a
\emph{thick arc}.  The photon's $\Delta = 0.245$ is a
\emph{thin arc} --- the most asymmetric in the catalog.  Every entry in
Table~\ref{tab:nuclear} has acquired geometric meaning:
helium-4's $\Delta = 0.064$ (thick arc) versus iron-56's
$\Delta = 0.036$ (thick arc approaching perfect ring).
The numbers have not changed.  What has changed is the reader's
capacity to \emph{read} them: the budget surface gave the gap
a geometry it did not have at the subatomic scale.  This is the
compounding at work --- each section deepens every previous one.

The diagnostic toolkit now has three registers:
(i)~the arithmetic of slaughter ($\Delta$ as a number),
(ii)~the dynamics of recovery (slaughter/recovery cycles at
phase boundaries), and (iii)~the geometry of the gap ($\Delta$
as lensing classifier, $\Gam'$ as gravity, $\Gam''$ as tidal
force).  Together they yield a further register:
\emph{prediction}.


\subsection{The Cosmic Web}

\begin{table}[htbp]
\caption{\label{tab:cosmic}Diffuse and composite structures.}
\begin{ruledtabular}
\begin{tabular}{llcccc}
Entity & Cat.\ & $F$ & $\IC$ & $\Delta$ & $\IC/F$ \\
\hline
Molecular cloud    & Diff.\ & 0.306 & 0.251 & 0.056 & 0.819 \\
Dust lane          & Diff.\ & 0.235 & 0.182 & 0.053 & 0.773 \\
Nebula             & Diff.\ & 0.354 & 0.305 & 0.048 & 0.863 \\
Intergal.\ medium  & Diff.\ & 0.174 & 0.112 & 0.062 & 0.643 \\
\hline
Spiral galaxy      & Comp.\ & 0.766 & 0.759 & 0.007 & 0.991 \\
Galaxy cluster     & Comp.\ & 0.706 & 0.700 & 0.006 & 0.991 \\
Binary star        & Comp.\ & 0.772 & 0.761 & 0.012 & 0.984 \\
Accretion disk     & Comp.\ & 0.591 & 0.581 & 0.010 & 0.983 \\
Cosmic filament    & Comp.\ & 0.509 & 0.490 & 0.019 & 0.963 \\
\end{tabular}
\end{ruledtabular}
\end{table}

The diffuse media sit in \regime{Collapse} regime with
low $F$ and $\omega > 0.30$ --- structurally incoherent
substrates from which stars condense.  Their role in the
scale ladder parallels the confinement cliff: they are
not failures but \emph{preconditions}, their low coherence
creating the gradient ($\dd\Gam/\dd\omega > 0$) that drives
gravitational condensation.

The \emph{nebula} ($\IC/F = 0.863$, $F = 0.354$) is the
most coherent diffuse entity, with moderate trajectory
closure ($c_5 = 0.35$) and enough internal structure ---
ionization fronts, shock boundaries --- to sustain local
collapse-return cycles.  The \emph{molecular cloud}
($\IC/F = 0.819$, $F = 0.306$) sits lower: it has enough
structure to fragment into clumps ($c_5 = 0.30$) but not
enough to stabilize --- collapse is inevitable, and from that
collapse, stars are born.  The \emph{dust lane}
($\IC/F = 0.773$, $F = 0.235$) is dimmer still: its
near-zero gradient ($c_3 = 0.03$) reflects the absence of
internal dynamics; it absorbs and obscures but does not
actively cycle.  The \emph{intergalactic medium}
($\IC/F = 0.643$, $F = 0.174$) is the most incoherent
entity in the astrophysical catalog.  Its channels are
uniformly depleted ($c_3 = 0.01$, $c_6 = 0.05$) --- the
residuum of cosmic evolution, the substrate that remains
after all condensation has occurred.
The four diffuse entities form a coherence gradient:
nebula~$>$~cloud~$>$~dust~$>$~IGM, tracing the
structural pathway from active collapse to inert residue.

The composite structures tell the opposite story ---
assembly restores coherence.  The \emph{spiral galaxy}
($\IC/F = 0.991$, $\Delta = 0.007$) achieves near-perfect
uniformity from billions of stellar oscillators: its
channel profile is remarkably flat ($c_0 = 0.82$,
$c_1 = 0.88$, $c_5 = 0.80$, $c_7 = 0.78$), with no single
channel below $0.55$.  The \emph{galaxy cluster}
($\IC/F = 0.991$, $\Delta = 0.006$) matches it through a
different mechanism: where the spiral galaxy's uniformity
comes from rotational averaging, the cluster's comes from
statistical depth --- hundreds of galaxies averaging their
channel profiles into the narrowest gap in the cosmic
catalog.
The \emph{binary star} ($\IC/F = 0.984$, $\Delta = 0.012$)
is a two-body composite whose trajectory closure
($c_5 = 0.90$) and return rate ($c_1 = 0.92$) reflect
mutual orbital periodicity --- the simplest case of
composite coherence exceeding individual components.
The \emph{accretion disk} ($\IC/F = 0.983$,
$\Delta = 0.010$) achieves comparable coherence through
a different route: its high gradient ($c_3 = 0.80$)
and deep well ($c_2 = 0.70$) compensate for a lower
return rate ($c_1 = 0.55$), reflecting matter that spirals
inward reliably but does not return.  It is a repetitive
orbital pattern with reliable but shallow return.
The \emph{cosmic filament} ($\IC/F = 0.963$,
$\Delta = 0.019$) is the least coherent composite, sitting
at the boundary between thick arc and perfect ring.  Its
moderate well depth ($c_2 = 0.55$) and broad distribution
($c_3 = 0.25$) reflect the large-scale structure of the
universe: coherent enough to channel matter along
gravitational corridors, but too diffuse for the tight
uniformity of a galaxy or cluster.


\subsection{The Budget Surface Itself}

\begin{theorem}[T-ST-9: Intrinsic Flatness]\label{thm:T9}
The Gaussian curvature of $z(\omega, C) = \Gam(\omega) +
\alpha\,C$ vanishes identically: $K = 0$.  The surface is
developable --- all structure comes from the embedding.
\end{theorem}

\begin{remark}
Since $z$ is linear in $C$, $\partial^2 z/\partial C^2 = 0$
and $K = 0$ follows algebraically: $K = (z_{xx}z_{yy} -
z_{xy}^2)/(1 + z_x^2 + z_y^2)^2$, and $z_{yy} = 0$,
$z_{xy} = 0$ force $K = 0$ identically.  This mirrors a key
finding from the 44 structural identities: the Bernoulli
manifold is flat in Fisher coordinates ($g_F(\theta) = 1$,
verified to $< 10^{-16}$~\cite{paulus2025umcp}).  Numerical
verification: $K$ evaluated on a $100 \times 100$ grid over
$\omega \in [0.01, 0.99]$, $C \in [0, 1]$ returns exactly~0
at all 10,000 points.  (3/3 sub-tests.)
\end{remark}

Table~\ref{tab:translation} lists qualitative correspondences
between GR concepts and budget-surface quantities.  These are
\emph{vocabulary mappings}, not derivations: the budget surface
has qualitative features (monotonic gradient, convexity, a pole)
that parallel features of gravitational potentials, but the
budget surface does not produce Einstein's field equations, and
the GR concepts listed have operational, experimentally confirmed
definitions that the budget surface does not replicate.

\begin{table}[htbp]
\caption{\label{tab:translation}Qualitative correspondences
between general relativity and GCD budget-surface quantities.
These are vocabulary mappings that highlight shared qualitative
features, not claims of physical equivalence.}
\begin{ruledtabular}
\begin{tabular}{lp{0.42\columnwidth}}
GR Concept & Budget-Surface Analog \\
\hline
Grav.\ binding & Accumulated $|\kappa|$ from iterations \\
Spatial curvature & Extrinsic curvature of $z(\omega, C)$ \\
Geodesic & Least-budget path (min-$\Gam$) \\
Cost divergence & $\omega \to 1$ pole ($\Gam \to \infty$) \\
Min.\ deposit & Minimum $|\Delta\kappa|$ per iteration \\
Cost near wells & $\tR$ increases near wells \\
Grad.\ steepening & $\dd^2\Gam/\dd\omega^2$ \\
Closed deflection & Closed deflection orbit \\
Deflect.\ asymmetry & Controlled by $\Delta = F - \IC$ \\
\end{tabular}
\end{ruledtabular}
\end{table}

As $\omega \to 1$, the denominator $1 - \omega + \eps \to \eps$
and $\Gam \to 1/\eps = 10^8$.  This pole means the cost of
return diverges --- a property of the rational function
$\omega^3/(1 - \omega + \eps)$, not a model of any physical
singularity.  The guard band $\eps$ prevents a true
divergence at the computational level.


% ============================================================
\section{Scale IV: The Living System}\label{sec:biology}
% ============================================================

The same kernel that classifies subatomic particles and
stellar objects now turns to the living body.  The Diagnostic
Key (Sec.~\ref{sec:key}) implies a structural expectation:
periodic oscillators with reliable return should have high
$\IC/F$ (the two-step pattern: periodicity provides the
degrees of freedom that close the gap); composites should
exceed their components (as nuclear binding exceeds hadron
coherence, Table~\ref{tab:nuclear}).

\textbf{What the Diagnostic Key implies.}
The reader can verify that the Diagnostic Key's structural
logic, applied to the encoding of Table~\ref{tab:channels},
implies the hierarchy before consulting
Table~\ref{tab:biological}.
The heart is a periodic oscillator with
reliable return --- the Key implies its $\IC/F$ should be
closer to the neutron star's $0.998$ (Table~\ref{tab:stellar},
Diagnostic Key item~7) than the dust lane's $0.773$
(Table~\ref{tab:cosmic}).  The cortex assembles from
individual neurons, as a galaxy cluster assembles from
individual stars --- the Key implies its $\IC/F$ should be
higher than the single neuron's.  The two-step pattern
implies ``closer to the neutron star'' and ``higher
than the single neuron'': periodic oscillators have high
$\IC/F$; composites exceed their components.
Note that this is a \emph{retroduction}: the encoding was
designed with knowledge of these systems, so the structural
expectation is a consistency check, not a blind prediction.
Let us check.

Four biological systems --- the neuron, the
cortical network, the cardiac rhythm, and the immune response
--- are encoded as 8-channel trace vectors using channels
that carry exactly the same operational meaning as at
subatomic and cosmic scales: coherence persistence, cycle
return rate, well depth, gradient strength, tidal symmetry,
trajectory closure, circulation area, and heterogeneity
profile.

\begin{table}[t]
\caption{\label{tab:biological}Biological entities.}
\begin{ruledtabular}
\begin{tabular}{lcccc}
Entity & $F$ & $\IC$ & $\Delta$ & $\IC/F$ \\
\hline
Neuron            & 0.679 & 0.633 & 0.046 & 0.933 \\
Cortical network  & 0.622 & 0.606 & 0.016 & 0.974 \\
Heart rhythm      & 0.769 & 0.730 & 0.039 & 0.949 \\
Immune response   & 0.588 & 0.575 & 0.013 & 0.978 \\
\end{tabular}
\end{ruledtabular}
\end{table}

\textbf{The neuron is a collapse-return oscillator.}
The action potential is a prototypical collapse-return cycle:
depolarization collapses the resting state, and the refractory
period guarantees return.  The neuron's $\IC/F = 0.933$ is
comparable to the carbon atom (0.898) and the stellar main
sequence (0.961) --- all three are moderately bound systems
whose internal heterogeneity produces a gap
$\Delta \approx 0.03$--$0.07$ without killing any channel
outright.
The weakest channel is gradient strength ($c_3 = 0.30$)
--- the single neuron's voltage gradient is modest ---
which introduces a gap of $\Delta = 0.046$.  But unlike
the proton or the black hole, no channel drops below $0.30$:
the neuron avoids geometric slaughter.  This is why neurons
can assemble into cortical networks that \emph{exceed} their
components (cortex $\IC/F = 0.974$ vs.\ neuron $0.933$) ---
the floor is high enough to survive composition.

\textbf{The heart is the fastest biological oscillator.}
The cardiac sinus rhythm achieves $\IC/F = 0.949$ and
$F = 0.769$ --- the highest fidelity among biological
entities.  Its trajectory closure ($c_5 = 0.95$) and
return rate ($c_1 = 0.98$) rival the pulsar's ($c_5 = 0.95$,
$c_1 = 0.98$).  The heart is a biological pulsar: both
the cardiac sinus node and the pulsar magnetosphere produce
rhythmic collapse-return cycles with millisecond-to-second
periods, and both achieve $\IC/F > 0.94$.
Its \regime{Watch} regime means it is the only biological
system stable enough to avoid \regime{Collapse}
classification.  When the heart's rhythm destabilizes
(arrhythmia), the structural signature is a drop in $c_5$
and $c_1$ --- the same channel degradation that separates
the neutron star from the black hole.

\textbf{The cortical network is a composite well.}
Like a galaxy cluster assembling from individual galaxies,
the cortex assembles from individual neurons.  Its
$\IC/F = 0.974$ is \emph{higher} than the single neuron's
$0.933$, reflecting the power of composite structure: the
network as a whole is more uniform than its parts, just as
the galaxy cluster ($\IC/F = 0.991$) exceeds the individual
stellar objects.  The structural reason is the same at both
scales: averaging over many oscillators smooths the channel
profile, raising the geometric mean toward the arithmetic
mean.  The gap shrinks from the neuron's $\Delta = 0.046$ to
the cortex's $\Delta = 0.016$ --- the same compression that
takes individual stars ($\Delta \approx 0.01$--$0.04$) to
the galaxy cluster's $\Delta = 0.006$.

\textbf{The immune response is a delayed-return cycle.}
The inflammatory-resolution cycle has $\IC/F = 0.978$ but
$F = 0.588$ --- moderate overall coherence.  The memory B
cells enable immunological return ($c_1 = 0.80$), but the
tidal symmetry ($c_4 = 0.60$) reflects the chronic
imbalance between pro- and anti-inflammatory signals.
The immune system is structurally in \regime{Collapse}
($\omega = 0.412 > 0.30$), not because it is failing, but
because its breadth of response keeps drift high.  This
mirrors the nebula ($\IC/F = 0.863$, $\omega > 0.30$) ---
both are diffuse, reactive systems operating in Collapse
regime as a functional requirement rather than a pathology.

\textbf{Biological time has a tempo hierarchy.}
The two-mechanism framework from Sec.~\ref{sec:time} applies
directly.  The heart ($c_1 = 0.98$, $\Gam' = 0.23$)
produces fast, cheap temporal ticks --- the highest
temporal resolution among biological entities.  The neuron ($c_1 =
0.90$) is slower: $10\%$ of action potentials fail
to produce a spike-response at the postsynaptic terminal,
and each tick costs roughly twice as much ($\Gam' = 0.53$).
The cortical network ($c_1 = 0.82$) is slower still:
network-level oscillations (alpha, gamma) emerge from
averaging many neuronal clocks, and the effective temporal
resolution coarsens --- this is why subjective time feels
smoother than neuronal time.  The immune system ($c_1 =
0.80$) operates on the slowest biological timescale: its
inflammatory-resolution cycles take days, and $20\%$ of
cycles fail to resolve completely.  This is why immune
recovery \emph{feels} slow, while heartbeats feel
instantaneous.  The hierarchy mirrors what we saw at the
astrophysical scale (pulsar $\to$ main sequence $\to$ red
giant $\to$ black hole), now re-instantiated at the
cellular level: same mechanism, different substrate, same
measurable ordering.

These biological well depths, gradients, and return rates
are computed by the same kernel, with the same frozen
parameters, using the same identities as those applied to
quarks and galaxies.  The mathematics does not change at the
phase boundary between physics and biology.  Only the channel
encoding changes --- and the structure that survives collapse
is measured identically.

\textbf{Consistency confirmed.}
The heart's $\IC/F = 0.949$ sits between the pulsar
($0.995$) and the neutron star ($0.998$): periodic
oscillators cluster together across twenty orders of
magnitude.  The cortical network's $\IC/F = 0.974$ exceeds
the single neuron's $0.933$, just as the galaxy cluster
($0.991$) exceeds individual stellar objects.  The structural
expectation matched the data.  This is a consistency
check, not a blind prediction --- the encoding was
designed with knowledge of these systems (see
Remark~\ref{rem:encoding}).  What is noteworthy is that
the \emph{same} structural logic (periodic oscillators
cluster high, composites exceed components) applies
across domains without domain-specific tuning.


% ============================================================
\section{Scale V: The Memory Well}\label{sec:memory}
% ============================================================

We arrive at the human mind.  The same 8-channel kernel
(Diagnostic Key, Sec.~\ref{sec:key}) now examines the most
intimate structures a human being possesses: memories,
habits, and grief.  If the two-step pattern
(geometric slaughter at phase boundaries, recovery through
new degrees of freedom) operates here, then trauma ---
which fragments specific channels while leaving others
vivid --- should produce the same diagnostic signature
as physical confinement.

\textbf{What the Diagnostic Key implies.}
The Diagnostic Key provides enough structure to derive
expectations --- noting that these are retroductions
(the encoding was chosen with clinical knowledge), not
blind predictions:
\begin{enumerate}
\item PTSD fragments specific channels while leaving others
  vivid.  This is heterogeneous --- expect $\Delta > 0.10$.
\item Avoidance prevents narrative completion --- expect
  trajectory closure $c_5 \approx 0.1$.
\item Intrusion-numbing asymmetry violates tidal symmetry ---
  expect $c_4 < 0.25$.
\item The gap size places it in thin-arc lensing
  (Theorem~\ref{thm:T7}).
\item The structural parallel is to confined quarks inside
  hadrons (the proton's $\IC/F = 0.718$,
  Table~\ref{tab:subatomic}): expect $\IC/F \approx
  0.72$--$0.76$.
\end{enumerate}
Compare this derivation against Table~\ref{tab:cognitive}.

\begin{table}[t]
\caption{\label{tab:cognitive}Cognitive entities.}
\begin{ruledtabular}
\begin{tabular}{lcccc}
Entity & $F$ & $\IC$ & $\Delta$ & $\IC/F$ \\
\hline
Healthy memory     & 0.781 & 0.767 & 0.014 & 0.982 \\
Trauma (PTSD)      & 0.481 & 0.366 & 0.115 & 0.760 \\
Habit loop         & 0.715 & 0.626 & 0.089 & 0.876 \\
Grief cycle        & 0.556 & 0.550 & 0.006 & 0.989 \\
\end{tabular}
\end{ruledtabular}
\end{table}


\subsection{Healthy Memory: The Perfect Ring}

A healthy memory has $F = 0.781$, $\IC = 0.767$, and
$\Delta = 0.014$.  All eight channels are active: high
coherence persistence ($c_0 = 0.85$), excellent recall return
($c_1 = 0.92$), moderate well depth ($c_2 = 0.65$), balanced
gradient ($c_3 = 0.55$), symmetric tidal forces
($c_4 = 0.90$), tight trajectory closure ($c_5 = 0.88$),
moderate circulation ($c_6 = 0.60$), and uniform channel
profile ($c_7 = 0.90$).  The lensing classification is
\emph{perfect ring}: the gap is too small to distort the
well's symmetry.

In the scale ladder, the healthy memory sits between the
stellar main sequence ($\IC/F = 0.961$) and the gas giant
($\IC/F = 0.953$), occupying a structurally familiar
position --- a deep, symmetric well with reliable return.
But its closest structural analog is the neutron star:
both achieve $\Delta < 0.015$ (\emph{perfect ring}), both
have near-uniform channel profiles, and both sit atop the
coherence peak of their respective recovery sequences.
The neutron star is the culmination of nuclear recovery
from confinement; the healthy memory is the culmination of
cognitive integration from neural substrate.  In both cases,
the coherence was built through iterated collapse-return
cycles across a phase boundary --- nuclear binding across
the hadron-nucleus boundary, and synaptic consolidation
across the neuron-network boundary.
The regime is \regime{Watch}: not perfectly stable (some
drift from perfect coherence), but far from collapse.


\subsection{PTSD: The Confinement Cliff of the Mind}

Now compare:
\beq\label{eq:ptsd_cliff}
  \frac{\IC}{F}\bigg|_{\text{healthy}} = 0.982
  \;\;\longrightarrow\;\;
  \frac{\IC}{F}\bigg|_{\text{PTSD}} = 0.760.
\eeq

This drop ($-22.6\%$) is structurally identical to the
confinement cliff:
\beq\label{eq:quark_cliff}
  \frac{\IC}{F}\bigg|_{\text{quark}} = 0.957
  \;\;\longrightarrow\;\;
  \frac{\IC}{F}\bigg|_{\text{proton}} = 0.718.
\eeq

The mechanism is the same.  The PTSD trace vector is:
\begin{center}
\texttt{(0.75, 0.30, 0.80, 0.70, 0.20, 0.10, 0.85, 0.15)}
\end{center}

Four channels are critically low:
\begin{itemize}
\item $c_1 = 0.30$: \textbf{cycle return rate} --- avoidance
  behavior means the memory is entered but not resolved.
  The mind falls into the well but does not return.
\item $c_4 = 0.20$: \textbf{tidal symmetry} ---
  hyperarousal and dissociation create a violently
  asymmetric well.  The pull is strong in one direction
  (intrusion) and absent in the other (numbing).
\item $c_5 = 0.10$: \textbf{trajectory closure} --- the
  re-experiencing cycle does not close.  The patient
  relives the trauma but cannot complete the narrative arc.
  This is the dead channel that kills IC.
\item $c_7 = 0.15$: \textbf{heterogeneity profile} ---
  extreme channel unevenness.  Some aspects of the memory
  are vivid ($c_2 = 0.80$, $c_6 = 0.85$); others are
  fragmentary.
\end{itemize}

The parallels to the proton are exact (Table~\ref{tab:parallel}).

\begin{table}[t]
\caption{\label{tab:parallel}Structural parallel: confinement
and PTSD.  The same geometric slaughter mechanism operates at
both scales.}
\begin{ruledtabular}
\begin{tabular}{lcc}
Feature & Proton & PTSD \\
\hline
$\IC/F$ & 0.718 & 0.760 \\
Dead channels & trajectory, het.\ & trajectory, het.\ \\
$\Delta$ & 0.134 & 0.115 \\
Regime & \regime{Collapse} & \regime{Collapse} \\
Lensing & thin arc & thin arc \\
What is confined & quarks & traumatic content \\
What cannot close & quark orbits & narrative arc \\
\end{tabular}
\end{ruledtabular}
\end{table}

\begin{remark}[Sensitivity of the Confinement Cliff]
\label{rem:sensitivity}
The PTSD classification as a ``confinement cliff'' depends on
the trajectory-closure channel value $c_5 = 0.10$.  If $c_5$
is perturbed:
\begin{itemize}
\item $c_5 = 0.10$ (baseline): $\IC/F = 0.760$, $\Delta = 0.115$
  --- confinement cliff, thin-arc lensing.
\item $c_5 = 0.30$: $\IC/F = 0.823$, $\Delta = 0.085$
  --- moderate depression, thick-arc lensing.
\item $c_5 = 0.50$: $\IC/F = 0.873$, $\Delta = 0.061$
  --- mild gap, thick-arc lensing.
\item $c_5 = 0.70$: $\IC/F = 0.916$, $\Delta = 0.041$
  --- approaching healthy range.
\end{itemize}
The cliff classification ($\IC/F < 0.80$, thin-arc) holds for
$c_5 \lesssim 0.20$.  Above $c_5 = 0.30$, the parallel to
the proton dissolves: the system transitions from thin arc
to thick arc, and the $\IC/F$ drop is a gentle slope
rather than a cliff.  The clinical literature on PTSD
supports $c_5 < 0.20$ (severe narrative fragmentation),
but the classification's encoding-dependence is real.
\end{remark}

The proton's quarks cannot close trajectories because
confinement prevents free propagation.  The PTSD patient's
memories cannot close trajectories because avoidance prevents
narrative completion.  In both cases, the dead trajectory
channel destroys the geometric mean while the arithmetic
mean (overall fidelity) remains
moderate~\cite{ehlers2000cognitive,vanderkolk2014body}.

\begin{remark}[Encoding Sensitivity of the Parallel]
\label{rem:parallel}
The proton--PTSD parallel (Table~\ref{tab:parallel})
depends on two encoding decisions: (i)~that ``trajectory
closure'' ($c_5$) is a meaningful channel for both quark
confinement and narrative completion, and (ii)~that both
receive low values ($c_5 = 0.15$ and $c_5 = 0.10$
respectively).  If $c_5$ were defined differently --- for
example, as spatial volume closure for quarks but temporal
narrative closure for memories --- the mapping would be a
\emph{conceptual analogy imposed by the encoder}, not a
structural discovery by the kernel.  Similarly, if the PTSD
trace vector were encoded with $c_5 = 0.50$ (moderate
narrative closure rather than severe), the $\IC/F$ drop
would be $0.982 \to 0.873$ --- a mild depression rather
than a confinement cliff.  The parallel is
\emph{encoding-contingent}: it is precise and reproducible
given the encoding in Table~\ref{tab:channels}, but
alternative encodings of the same clinical phenomena could
produce different structural classifications.  What this
framework offers is a \emph{controlled comparison}: once
the encoding is fixed, the structural diagnostics are
determined, and the question becomes whether the same
encoding scheme, applied consistently across domains,
produces parallels that domain experts recognize as
meaningful.
\end{remark}

The lensing classification confirms the analogy: both the
proton and PTSD are classified as \emph{thin arc} ---
moderate asymmetry distorting the well gradient into a
one-sided pull.  A healthy memory, like a neutron star, is
a \emph{perfect ring}.

\textbf{PTSD freezes time.}
The two-mechanism framework explains why.  With $c_1 = 0.30$,
only $30\%$ of collapse-return cycles complete --- $70\%$ of
temporal ticks are lost.  Simultaneously, $\Gam'(0.519) =
2.28$ makes each surviving tick $11\times$ more expensive than
a healthy memory's.  The effective temporal resolution
$c_1/\Gam' = 0.30/2.28 = 0.13$ is the lowest of any
entity in this paper --- worse even than the black hole
($c_1/\Gam' = 0.50/0.48 = 1.04$), because the trauma well
combines low return rate with a steep budget gradient.  This is the structural substrate of what
clinicians call \emph{temporal disintegration}: PTSD patients
report that the traumatic event feels as though it is
happening ``right now,'' that time between intrusions
collapses, that minutes in a trigger state feel like hours.
In the diagnostic framework, this corresponds to an
accurate reading of the system's temporal structure:
time, as measured by $\mathcal{T} = c_1/\Gam'$,
moves differently in a fragmented well.

\textbf{Clinical consequence as diagnostic re-description.}
The kernel identifies which channels are low:
restore $c_5$ (trajectory closure) through
narrative completion --- prolonged exposure therapy;
restore $c_4$ (tidal symmetry) through grounding techniques
that equilibrate intrusion and numbing; restore $c_1$
(return rate) through graduated re-entry and
desensitization.  However, this is a re-description of
clinical knowledge already encoded in the channel values,
not a derivation from the mathematics: avoidance was
encoded as low $c_5$, so the kernel's recommendation to
``restore $c_5$'' restates the encoding.
What the kernel adds is the \emph{budget context}: the
gradient $\Gam'(\omega)$ at the trauma's $\omega = 0.519$
is $11\times$ stronger than at the healthy memory's
$\omega = 0.219$, suggesting that recovery should proceed
by small steps down the budget surface, closing channels
one at a time rather than all at once --- a structural
rationale for graduated exposure that parallels how
nuclear binding replaces dead hadron channels incrementally
through shell filling and pairing.

The diagnostic toolkit has accumulated its final register:
arithmetic (the gap), dynamics (slaughter and recovery),
geometry (lensing and gravity), consistency (biology
confirmed structural expectations), and now
\emph{re-description} --- the kernel identifies which
channels are dead and provides a cost-based rationale
for the order of restoration.  Each register
enriches all the others.

\textbf{The black hole parallel.}
The black hole ($\IC/F = 0.825$) has dead trajectory closure
($c_5 = 0.10$) and reduced return rate ($c_1 = 0.50$) ---
the same two channels that are near-dead in PTSD ($c_5 = 0.10$,
$c_1 = 0.30$).  The clinical observation that PTSD creates
an ``inescapable pull'' toward re-experiencing has a
diagnostic counterpart in this framework:
the budget gradient $\dd\Gam/\dd\omega$ at the trauma well's
$\omega = 0.519$ produces a pull of $\Gam'(0.519) = 2.28$,
while the healthy memory at $\omega = 0.219$ experiences
$\Gam'(0.219) = 0.20$ --- an $11\times$ stronger gravitational
pull.

As $\omega \to 1$, the cost of return diverges:
$\Gam \to 1/\eps = 10^8$.  This is the event horizon of the
budget surface.  PTSD is a system at $\omega = 0.519$, well
within the \regime{Collapse} regime ($\omega \geq 0.30$),
experiencing significant but not infinite cost of return.
The therapeutic implication is structural: recovery requires
restoring the dead channels --- closing the trajectory
($c_5$), balancing the tidal symmetry ($c_4$), improving
the return rate ($c_1$) --- exactly as nuclear binding
restores the channels that confinement killed.


\subsection{The Habit Loop: A Shallow Symmetric Well}

The habit loop ($F = 0.715$, $\IC = 0.626$,
$\Delta = 0.089$) is the structural opposite of PTSD.
Its trajectory closure ($c_5 = 0.92$) and return rate
($c_1 = 0.95$) are the highest among cognitive entities.
The well is shallow (moderate depth, low gradient) and
symmetric (high tidal symmetry).  The regime is
\regime{Watch}: the habit is stable enough to avoid
collapse, maintained by reliable cycle return.

The habit loop's closest astrophysical analog is the rocky
planet ($\IC/F = 0.875$, $\Delta = 0.082$): both are
shallow wells maintained by repetitive orbital dynamics
rather than deep binding.  The planet orbits its star;
the habit loop orbits its cue-routine-reward cycle.  Neither
has the depth of the neutron star or the healthy memory,
but both have the trajectory closure that the black hole
and PTSD lack.  The accretion disk ($\IC/F = 0.983$)
offers another parallel: a repetitive structure where
material circulates reliably, but the well is defined by
the orbit rather than by internal binding.

Where PTSD is a black hole --- deep, asymmetric, hard to
escape --- the habit loop is a gentle orbit, a well-worn
track in the budget surface.  Temporally, the habit loop is
the cognitive inverse of PTSD: with $c_1 = 0.95$ and
$\Gam'(0.285) = 0.39$, its effective temporal resolution
$c_1/\Gam' = 2.46$ is the highest among cognitive entities.
Time in a habit loop moves \emph{fast} --- cycles complete
reliably and cheaply, producing the subjective experience
of time vanishing: hours pass in routine without being
noticed.  This is the cognitive analog of the pulsar
($c_1/\Gam' \approx 14$), where temporal ticks are so
numerous and so cheap that each individual tick carries
almost no experiential weight.  The clinical correlate
is automatic behavior: habits feel timeless precisely
because their temporal resolution is so high that
individual moments blur together.  The lensing classification
is \emph{thick arc}: some heterogeneity
($\Delta = 0.089$) from the depth and gradient channels,
but no channel is dead enough to produce geometric
slaughter.  This is why habits are easy to break
compared to trauma: the budget gradient at the habit's
$\omega = 0.285$ is $\Gam'(0.285) = 0.39$, six times
weaker than at the PTSD well's $\omega = 0.519$.  The pull
is real but escapable.


\subsection{Grief: Uniform Depression}

Grief ($F = 0.556$, $\IC = 0.550$, $\Delta = 0.006$) is
structurally distinct from PTSD.  Its $\IC/F = 0.989$ is
the \emph{highest} among cognitive entities --- the channels
are uniformly depressed rather than selectively killed.
The lensing classification is \emph{perfect ring}: the well
is symmetric, with no dominant dead channel.

The grief cycle's closest astrophysical analog is not
the black hole --- it is the dust lane ($\IC/F = 0.773$,
$\Delta = 0.053$) or the molecular cloud ($\IC/F = 0.819$,
$\Delta = 0.056$).  Like diffuse astrophysical media, grief
is characterized by uniformly low coherence rather than
selective channel death.  Everything is reduced;
nothing is shattered.  The dust lane is dim across all
wavelengths; the grieving mind is diminished across all functions.
But crucially, the $\IC/F$ remains high because the
reduction is symmetric --- the geometric mean tracks the
arithmetic mean closely.  This is why, unlike PTSD, grief
does not produce geometric slaughter: there is no single
dead channel dragging the geometric mean to the floor.

This captures a clinical truth.  In grief, everything hurts
--- all channels are moderately reduced.  In PTSD, specific
channels are shattered while others remain vivid.  The kernel
distinguishes these two states precisely: grief has low $F$
(low overall coherence) but high $\IC/F$ (uniform channels);
PTSD has moderate $F$ but low $\IC/F$ (fragmented channels).
The heterogeneity gap $\Delta$ tells the difference:
$\Delta_{\text{grief}} = 0.006$,
$\Delta_{\text{PTSD}} = 0.115$ --- a $19\times$ ratio.
Grief's $\IC/F$ (0.989) exceeds the healthy memory's (0.982)
--- a structural paradox that resolves once one recognizes
that $\IC/F$ measures channel \emph{uniformity}, not
well-being: grief is uniformly diminished, and uniformity
is what the geometric mean rewards.

Temporally, grief occupies the middle ground.  With $c_1 =
0.55$ and $\Gam'(0.444) = 1.34$, the effective temporal
resolution is $c_1/\Gam' = 0.41$ --- slower than the healthy
memory ($4.58$) and the habit loop ($2.46$), but faster than
PTSD ($0.13$).  Grief does not freeze time; it \emph{slows}
it.  The grieving person experiences days that feel like
weeks, not because cycles fail to complete (as in PTSD), but
because each cycle costs more than it should.  The uniform
channel depression means that every domain of experience ---
work, rest, social interaction --- is equally sluggish.
There is no fragmentation, no dead channel, no frozen moment;
there is only a pervasive temporal drag, as though the entire
system is wading through the budget surface at a higher
altitude than normal.

The therapeutic implication follows from the structure.
PTSD recovery requires restoring specific dead channels
(trajectory closure, tidal symmetry) --- targeted
intervention, like nuclear binding targeting specific
dead hadron channels.  Grief recovery requires raising all
channels simultaneously --- a general restoration of
coherence, more akin to how a molecular cloud condenses
into a star: the whole substrate must reach sufficient
density before structure can nucleate.


% ============================================================
\section{The Unification}\label{sec:unification}
% ============================================================

\begin{remark}[Theorem Classification]
\label{rem:theorem_classes}
The ten theorems fall into two classes:
\begin{itemize}
\item \textbf{Mathematical} (properties of $\Gam(\omega) =
  \omega^3/(1 - \omega + \eps)$, holding for any input):
  T-ST-1 (gradient positive), T-ST-2 (gradient increasing),
  T-ST-3 (cubic onset), T-ST-5 (gradient depends on $\omega$
  alone), T-ST-9 (intrinsic flatness $K = 0$).
\item \textbf{Encoding-dependent} (structural observations
  that depend on the channel values assigned to the
  40~entities): T-ST-4 (time dilation), T-ST-6 (memory wells),
  T-ST-7 (lensing classification), T-ST-8 (arrow of time),
  T-ST-10 (cross-scale consistency).
\end{itemize}
The mathematical theorems verify calculus properties of a
rational function.  The encoding-dependent theorems document
patterns that emerge from the specific channel assignments
in the entity catalog.  Different encodings of the same
phenomena could produce different patterns
(Remark~\ref{rem:encoding}, Remark~\ref{rem:sensitivity}).
\end{remark}

\begin{theorem}[T-ST-10: Cross-Scale Consistency]
\label{thm:T10}
All 40~entities in all 9~categories obey the three Tier-1
kernel identities:
\begin{enumerate}
\item $F + \omega = 1$ to $< 10^{-12}$ for all entities.
\item $\IC \leq F$ for all entities.
\item $\IC = e^{\kappa}$ to $< 10^{-10}$ for all entities.
\end{enumerate}
The trauma well has lower $\IC$ than the healthy memory
($\IC_{\text{PTSD}} = 0.366 < \IC_{\text{healthy}} = 0.767$).
All 9~categories are represented and non-empty.
\end{theorem}

\begin{remark}
The consistency is not a claim that quarks and memories
``are the same.''  It is the observation that a consistent
encoding scheme, when processed through the kernel, produces
the same diagnostic pattern at every scale: geometric
slaughter at phase boundaries followed by recovery through
new degrees of freedom.  The three algebraic identities are
mathematical constraints that hold for \emph{any} input in
$[0,1]^n$ --- they must hold, by construction (see
Remark~\ref{rem:theorem_classes}).  What is
empirical is the encoding: the channel values assigned from
domain-specific observables.  The significant finding is not
that the identities hold (they must) but that the
\emph{diagnostic values} ($F$, $\IC$, $\Delta$, $\mathcal{T}$)
cluster into the same structural types (geometric slaughter
with $\IC/F \approx 0.7$, perfect rings with $\Delta < 0.02$,
cost-frozen temporal resolution) across systems separated by
twenty orders of magnitude.  Whether this clustering reflects
a genuine structural motif of nature or a systematic feature
of how humans encode oscillatory systems remains an open
question --- but the pattern, once exhibited, is precise
enough to be falsified (Remark~\ref{rem:encoding}).
(6/6 sub-tests.)
\end{remark}

The complete scale ladder is summarized in
Table~\ref{tab:ladder}.

\clearpage
\begin{table*}[t]
\caption{\label{tab:ladder}The scale ladder: geometric
slaughter and recovery across twenty orders of magnitude.
$\langle \IC/F \rangle$ is the category mean.}
\begin{ruledtabular}
\begin{tabular}{llcccl}
Scale & Category & Entities & $\langle F \rangle$ &
  $\langle \IC/F \rangle$ & Structural Signature \\
\hline
$10^{-15}$\,m & Subatomic &
  5 & 0.509 & 0.656 &
  Confinement cliff: quark 0.957 $\to$ proton 0.718 \\
$10^{-14}$\,m & Nuclear/Atomic &
  3 & 0.747 & 0.923 &
  Recovery: binding restores dead channels \\
$10^{9}$--$10^{4}$\,m & Stellar &
  7 & 0.795 & 0.957 &
  Deep symmetric wells; neutron star 0.998 \\
$10^{6}$--$10^{8}$\,m & Planetary &
  5 & 0.632 & 0.856 &
  Shallower wells; Moon lowest at 0.742 \\
$10^{16}$--$10^{23}$\,m & Diffuse &
  4 & 0.267 & 0.775 &
  Low $F$, high $\omega$ --- precursors to structure \\
$10^{20}$--$10^{25}$\,m & Composite &
  5 & 0.669 & 0.982 &
  Assembled wells --- highest category $\IC/F$ \\
$10^{-6}$--$10^{-3}$\,m & Biological &
  4 & 0.665 & 0.959 &
  Living oscillators: heart rivals the pulsar \\
--- & Cognitive &
  4 & 0.633 & 0.902 &
  Memory wells: PTSD = confinement of the mind \\
--- & Boundary &
  3 & 0.389 & 0.651 &
  No recovery: glass, turbulence, decoherence \\
\end{tabular}
\end{ruledtabular}
\end{table*}

At every phase boundary --- quarks to hadrons, hadrons to
nuclei, diffuse gas to stars, substrate to organism, neurons
to organized memory --- the same two-step process operates:

\begin{enumerate}
\item \textbf{Geometric slaughter at the boundary.}
  Dead channels in the trace vector destroy the geometric
  mean $\IC$ while the arithmetic mean $F$ survives.  The
  heterogeneity gap $\Delta$ opens.
\item \textbf{Recovery through new degrees of freedom.}
  New structure (nuclear binding, stellar fusion, synaptic
  integration, narrative completion) replaces dead channels,
  closing the gap and restoring $\IC/F$ toward unity.
\end{enumerate}

This pattern is the content of Axiom-0: \emph{collapse is
generative}.  The confinement cliff is not a catastrophe;
it is the precondition for nuclear binding.  The PTSD well
is not a permanent state; it is the precondition for
therapeutic recovery.  \emph{Only what returns is real.}

\begin{remark}[Entity Selection and Boundary Cases]
\label{rem:selection}
The 40~entities span the scale ladder from $10^{-15}$\,m to
macroscopic scales, with at least three entities per
category.  This is a \emph{curated catalog}, not a
representative sample.  To test the pattern's limits, three
boundary cases --- systems where the recovery step may fail
--- are included as a dedicated category.

\emph{Glass (amorphous solid)}: frozen disorder with no
periodic collapse-return cycles ($c_1 = 0.08$, $c_5 = 0.10$).
Result: $F = 0.405$, $\IC/F = 0.613$, Collapse regime.
Four dead channels create a heterogeneity gap ($\Delta =
0.157$) comparable to the subatomic mean --- but unlike
nuclei, there is no recovery step.  The glass \emph{stays}
in Collapse because the frozen structure never returns.

\emph{Turbulent flow}: energy cascades and vortex
circulation ($c_3 = 0.85$, $c_6 = 0.80$) but no trajectory
closure ($c_5 = 0.05$).  Result: $F = 0.513$, $\IC/F =
0.785$, Collapse regime.  This is the highest $\IC/F$ of the
three boundary cases --- turbulence is dynamically rich but
structurally incomplete, and the single dead closure channel
is partially compensated by strong gradient and circulation.

\emph{Decoherence event}: irreversible coherence loss
($c_1 = 0.02$, $c_5 = 0.03$, $c_6 = 0.05$).  Result:
$F = 0.250$, $\IC/F = 0.554$, deep Collapse ($\omega =
0.75$).  Remark~5 in the prior version predicted this would
occupy ``the same structural position as the photon.''  The
kernel reveals a more nuanced picture: both share $c_1
\approx 0$ (no return) and both are in Collapse, but the
photon's $\IC/F = 0.067$ reflects \emph{extremely} dead
channels ($c_2, c_3, c_6 < 0.001$), while the decoherence
event's channels, though low, are not as extreme.  The
prediction is partially confirmed (both are
no-return systems in Collapse) and partially refined (they
differ by an order of magnitude in $\IC/F$).

All three boundary cases satisfy the Tier-1 identities
to machine precision.  Their structural significance: they
are systems where geometric slaughter occurs but recovery
\emph{does not follow}.  This does not contradict the
two-step pattern; it confirms that the pattern has a
testable boundary --- recovery is not guaranteed, and the
kernel correctly classifies non-returning systems as
structurally incomplete.
\end{remark}

\textbf{Compounding as structural enrichment.}
The fact that understanding deepens through this paper is
not a rhetorical device --- it is a structural feature of
single-axiom systems.  Because every identity, every theorem,
and every classification derives from the same source
(Axiom-0 and the budget $\Gam$), each new derivation enriches
every previous one.  The gap $\Delta$ is richer as a lensing
classifier (Section~\ref{sec:cosmos}) than as a bare number
(Section~\ref{sec:subatomic}), and the clinical protocol
(Section~\ref{sec:memory}) is richer after one has seen the
recovery mechanism at nuclear scales.

But each section is also independently coherent: the
Diagnostic Key (Sec.~\ref{sec:key}) provides the fixed
toolkit, and each section makes its own case against it.
Just as the kernel reads any domain through the same
invariants, each section can be read through the same
Diagnostic Key.  The compounding is real --- each section
deepens all others --- but it is an enrichment, not a
prerequisite.  This is the \emph{geometry of
compounding}: in a system with a single axiom, every
derivation strengthens the whole, yet each derivation
stands on the axiom alone.


% ============================================================
\section{Reproducibility}\label{sec:howto}
% ============================================================

Every theorem in this paper is machine-verified.  The
reference implementation~\cite{umcpmetadatarepo} includes
the full entity catalog (40~entities, 9~categories), the
kernel computation, and the theorem-proving harness.  To
reproduce all results:

\begin{small}
\begin{verbatim}
pip install -e ".[all]"
python -c "
from closures.spacetime_memory\
  .spacetime_theorems \
  import prove_all_theorems
results = prove_all_theorems()
for r in results:
    print(f'{r.name}: {r.verdict}'
          f' ({r.n_passed}/{r.n_tests})')
"
\end{verbatim}
\end{small}

Output: 10/10~PROVEN, 37/37 sub-tests, 0~failures.  The
full test suite (\texttt{pytest
tests/test\_253\_spacetime\_memory\_theorems.py -v})
runs 175~tests in $< 300$~ms with zero failures.

The theorem suite is structured as follows.  Each theorem
object specifies its name, statement, and a list of
sub-tests.  Sub-tests are predicate functions that return
\texttt{True} or \texttt{False}; a theorem is PROVEN only
if all sub-tests pass.  The sub-test counts reported in
the Remarks throughout this paper (e.g., ``4/4 sub-tests''
for Theorem~\ref{thm:T1}) correspond exactly to the
programmatic sub-tests in the harness.  No theorem is
accepted on the basis of spot-checks or manual verification.

\textbf{Test count breakdown.}
Of the 175~tests, 37 are theorem sub-tests (the direct
verifications reported in the Remarks above).  The remaining
138 are infrastructure tests: entity completeness, channel
bounds, schema validation, category coverage, and
reproducibility checks.  These validate the implementation,
not the physics claims --- they ensure that the catalog is
correctly formed so that the theorem sub-tests can be
trusted.

The entity catalog is defined in
\texttt{spacetime\_kernel.py}, where each entity's 8-channel
trace vector is frozen as a Python tuple.  Modifying any
channel value and re-running the theorem suite will produce
different $F$, $\IC$, and $\Delta$ values --- and may cause
theorems to fail, demonstrating the encoding-sensitivity
discussed in Remark~\ref{rem:encoding}.  This is by design:
the theorems test the \emph{current} encoding, not a
universal truth.


% ============================================================
\section{Conclusions}\label{sec:conclusions}
% ============================================================

We have traced the geometry of what returns from quarks to
memory wells, using a single axiom, a single kernel, and a
single budget surface.  Across 40~entities and 9~categories
spanning twenty orders of magnitude, one structural pattern
recurs: \emph{structure accumulates through iterated
collapse-return cycles}, and the kernel measures how that
accumulation proceeds.  The pattern has three modes:

\begin{itemize}
\item \textbf{Space accumulates as iterated well depth.}
  Mass, in this diagnostic framework, is accumulated
  $|\kappa|$ from repeated collapse-return cycles
  (Theorem~\ref{thm:T6}).  A proton's well depth is the
  accumulated confinement of three quarks.  A neutron star's
  depth is the accumulated compression of $10^{57}$
  nucleons.  A trauma well's depth is the accumulated weight
  of repeated re-experiencing.  The budget gradient
  ($\dd\Gam/\dd\omega > 0$, Theorem~\ref{thm:T1}) ensures
  that deeper wells cost more to escape --- a diagnostic
  analog of what general relativity describes as
  gravitational binding.

\item \textbf{Time accumulates as affordable ticks.}
  A system's internal time runs at the rate at which it
  can complete and afford collapse-return cycles
  ($\mathcal{T} = c_1/\Gam'$, Definition~\ref{def:temporal}).
  The neutron star ticks cheaply and reliably
  ($\mathcal{T} \approx 18$); the PTSD well ticks rarely
  and expensively ($\mathcal{T} = 0.13$); the photon barely
  ticks at all ($\mathcal{T} \approx 0$).  The cost asymmetry
  (descent cheap, ascent expensive, Theorem~\ref{thm:T8})
  produces a directional arrow: in this diagnostic
  framework, time runs forward because return costs more
  than collapse.

\item \textbf{Memory accumulates as well geometry.}
  The shape of a well --- its channel profile, its gap
  $\Delta$, its lensing classification --- is the accumulated
  record of how collapse-return cycles have proceeded.  A
  healthy memory is a perfect ring ($\Delta = 0.014$): all
  channels uniformly reinforced by repeated successful
  recall.  PTSD is a thin arc ($\Delta = 0.115$): selective
  channels dead from avoidance, others vivid from intrusion.
  Grief is a uniform depression ($\Delta = 0.006$): all
  channels lowered, none shattered.  The well geometry
  encodes the memory --- the structural residue of every
  cycle that built it, rendered as a diagnostic profile.
\end{itemize}

The three modes are not independent.  Spatial depth,
temporal resolution, and well geometry are three readings of
the same underlying process: iterated collapse-return.  In
this diagnostic reading, mass corresponds to the depth the
iterations have carved; time, to the rate at which iterations
proceed; memory, to the shape the iterations have
left.  At every phase boundary --- quarks to hadrons,
diffuse gas to stars, neurons to organized memory --- the
two-step pattern operates: geometric slaughter destroys
coherence at the boundary, and new degrees of freedom
restore it on the other side.

\textbf{What this paper has shown.}
The introduction made five claims --- three structural, two
mathematical.  We close the loop:

\emph{Structural claims} (encoding-dependent; falsifiable
given a fixed encoding scheme):

\begin{enumerate}
\item \textbf{Space accumulates as iterated $|\kappa|$}
  (Theorem~\ref{thm:T6}).
  Linearity verified: $\mathrm{depth}(2N) = 2\,\mathrm{depth}(N)$
  to $< 10^{-12}$ for all entities.
  \emph{Falsified by}: an entity, under the same encoding,
  whose well depth does not grow linearly with iteration
  count.  (4/4 sub-tests.)

\item \textbf{Time emerges from return cost}
  (Definition~\ref{def:temporal}, Theorems~\ref{thm:T4},
  \ref{thm:T8}).
  $\mathcal{T} = c_1/\Gam'$ ranks all 40~entities from
  photon ($\approx 0$) to neutron star ($\approx 18$).
  Cost ratio: $\Gam(0.90)/\Gam(0.10) = 6.56 \times 10^{3}$.
  \emph{Falsified by}: an entity whose independently
  established temporal ordering contradicts $\mathcal{T}$
  (Remark~\ref{rem:temporal_choice}).
  (3/3 + 3/3 sub-tests.)

\item \textbf{Memory accumulates as well geometry}
  (Definition~\ref{def:slaughter},
  Theorems~\ref{thm:T6}, \ref{thm:T10}).
  Confinement cliff: $\IC/F$ drops $0.957 \to 0.718$ across
  color confinement; trauma cliff: $0.982 \to 0.760$ across
  narrative closure.  Same dead channels ($c_5$, $c_7$), same
  mechanism (Remark~\ref{rem:parallel}).
  \emph{Falsified by}: a phase boundary, under the same
  encoding, where the $\IC/F$ drop is not explained by dead
  channels.
\end{enumerate}

\emph{Mathematical claims} (hold for any valid input;
verify implementation correctness):

\begin{enumerate}\setcounter{enumi}{3}
\item \textbf{Budget gradient structures the manifold}
  (Theorem~\ref{thm:T1}).
  $\dd\Gam/\dd\omega > 0$ at 100~grid points and all
  40~positions; minimum $\Gam'(0.01) = 3.0 \times 10^{-4}$.
  This is a mathematical property of $\Gam$ with $p = 3$.
  (4/4 sub-tests.)

\item \textbf{Structural consistency across all scales}
  (Theorem~\ref{thm:T10}).
  $F + \omega = 1$ to $< 10^{-12}$; $\IC \leq F$ with zero
  violations; $\IC = e^{\kappa}$ to $< 10^{-10}$.  All
  9~categories represented.  These identities hold for any
  input in $[0,1]^n$; they verify the implementation, not
  the encoding.  (6/6 sub-tests.)
\end{enumerate}

\textbf{The structural parallel.}
The parallel between the proton and PTSD
(Table~\ref{tab:parallel}) is the paper's most striking
observation: both exhibit $\IC/F \approx 0.7$, both have dead
trajectory closure and heterogeneity channels, both are
classified as \regime{Collapse} with thin-arc lensing.  The
proton's quarks are confined because they cannot propagate
freely; the PTSD patient's trauma is confined because the
narrative cannot close.  The diagnostic mathematics is the
same.  In diagnostic terms, the recovery pathway is the
same: restore dead channels, close the trajectory, reduce
the gap.

This parallel depends on the encoding
(Remark~\ref{rem:encoding}): the channel values were
assigned from domain-specific observables, and different
assignments would produce different diagnostics.  What is
noteworthy is that when domain experts encode confinement
physics and clinical trauma using the same 8-channel schema,
the kernel reports the same structural signature.  Whether
this reflects a deep structural isomorphism between physical
and psychological confinement, or a systematic feature of how
humans conceptualize constraint, is a question the framework
makes precise enough to investigate.

\textbf{What the budget surface is not.}
No new physics is introduced.  The budget surface does not
replace general relativity or neuroscience.  It provides a
\emph{diagnostic lens} --- domain-agnostic, computationally
verifiable, and continuous-integration--ready --- through which
phenomena separated by twenty orders of magnitude in scale
become structurally visible as instances of a single
accumulation pattern.  The diagnostic mappings to GR concepts
(Table~\ref{tab:translation}) are structural correspondences,
not derivations --- they reveal that the budget surface has
the same qualitative features (monotonic gradient,
convexity, pole at $\omega \to 1$) as a gravitational
potential, without claiming to produce Einstein's field
equations.

\textbf{The triple accumulation.}
The reader who began this paper with a question --- ``What
does a quark have in common with a traumatic memory?'' ---
now holds the answer as a pattern: both are systems where
collapse-return cycles have accumulated into a well whose
geometry the kernel can measure.  The quark's well is
spatial: accumulated $|\kappa|$ from color confinement.  The
star's well is temporal: its internal clock ticks at a rate
set by the budget gradient.  The memory's well is geometric:
its shape records every recall cycle that carved it.  Space,
time, and memory exhibit parallel diagnostic signatures of
the same accumulation pattern --- three ways that iterated
collapse-return leaves a structural residue that the kernel
reads as $F$, $\IC$, and $\Delta$.

That progression of understanding --- from arithmetic, to
dynamics, to geometry, to consistency checks, to clinical
re-description --- is the enrichment that linear reading adds.
But the paper can also be entered at any section: the
Diagnostic Key (Sec.~\ref{sec:key}) is the fixed lens
through which every scale passes, and the same two-step
pattern emerges regardless of entry point.  The paper enacts
the geometry it describes.  Understanding that compounds is
understanding that returns --- and understanding that can
be entered from any angle is understanding whose coherence
is structural, not sequential.

\emph{Collapsus generativus est; solum quod redit, reale est.}

% ============================================================
\begin{acknowledgments}
The cross-scale insight --- that confinement and PTSD share
geometric-slaughter signatures --- emerged from the
orientation script's derivation chains.  Reference
implementation:
\url{https://github.com/calebpruett927/GENERATIVE-COLLAPSE-DYNAMICS}.
\end{acknowledgments}

\bibliography{Bibliography}

\end{document}
